% Jedan izvorni slajd = jedan frame = jedna PDF strana.
\documentclass[aspectratio=169,11pt]{beamer}
\newcommand{\ETFTemaPutanja}{../_zajednicko/etf-v1}
\makeatletter
\edef\input@path{{\ETFTemaPutanja/}}
\makeatother
\usepackage{stil,dijagrami}
\definecolor{DEcrvena}{HTML}{A32638}
\setbeamercolor{alerted text}{fg=DEcrvena}
\renewcommand{\gusttekst}{\fontsize{11}{13}\selectfont}
\title[Digitalna elektronika 1 — 2021/22]{Logička kola sa MOS tranzistorima — 1. deo}
\author{prof. dr Lazar Saranovac}
\date{2021/22}
\setlength{\abovedisplayskip}{5pt}
\setlength{\belowdisplayskip}{5pt}
\setlength{\abovedisplayshortskip}{3pt}
\setlength{\belowdisplayshortskip}{3pt}
\begin{document}
\slajdid{09_1-s001}
\begin{frame}[label=09_1-s001]{}\setlength{\parskip}{2pt}\setlength{\abovedisplayskip}{4pt}\setlength{\belowdisplayskip}{4pt}\setlength{\abovedisplayshortskip}{3pt}\setlength{\belowdisplayshortskip}{3pt}
% element: 09_1-s001:e01
% element: 09_1-s001:e02
\centering
{\fontsize{25}{29}\selectfont\bfseries Logička kola\par sa MOS tranzistorima\par}
\vspace{7mm}{\Large 1. deo\par}
\vspace{8mm}prof. dr Lazar Saranovac\par Katedra za elektroniku\par Digitalna elektronika 1 — 2021/22
\end{frame}

\slajdid{09_1-s002}
\begin{frame}[label=09_1-s002]{Simboli MOS tranzistora}\setlength{\parskip}{2pt}\setlength{\abovedisplayskip}{4pt}\setlength{\belowdisplayskip}{4pt}\setlength{\abovedisplayshortskip}{3pt}\setlength{\belowdisplayshortskip}{3pt}
% element: 09_1-s002:e01
% element: 09_1-s002:e02
% element: 09_1-s002:e03
\centering \begin{tikzpicture}[x=1cm,y=.48cm,font=\sitneoznake,>=Latex,line width=.7pt]
\foreach\col/\pol/\dep/\name in{0/0/0/nMOS,1/1/0/pMOS,2/0/1/nMOS,3/1/1/pMOS}{
 \node at(\col*2.7,3.9){\name};
 \foreach\row in{1,2,3,4}{\begin{scope}[shift={(\col*2.7,3.55-\row*2.2)}]
 \ifnum\row<3
  \draw(-.85,-.22)--(-.08,-.22);\draw(-.08,-.37)--(-.08,.37);
  \ifnum\dep=0 \foreach\lo/\hi in{-.4/-.22,-.1/.1,.22/.4}{\draw(0,\lo)--(0,\hi);}\else\draw[line width=1.1pt](0,-.4)--(0,.4);\fi
  \draw(0,.3)--(.4,.3)--(.4,.73)node[right]{D};\draw(0,-.3)--(.4,-.3)--(.4,-.73)node[right]{S};
  \ifnum\pol=0\draw[-{Triangle[length=2.5mm,width=1.8mm]}](.4,0)--(.04,0);\else\draw[-{Triangle[length=2.5mm,width=1.8mm]}](.04,0)--(.4,0);\fi
  \ifnum\row=1\draw(.4,0)--(.85,0)node[right]{B};\else\draw(.4,0)--(.4,-.3);\fi
  \node[above]at(-.77,-.22){G};
 \else
  \draw(-.08,-.32)--(-.08,.32);
  \ifnum\dep=1\draw[line width=1.3pt](0,-.4)--(0,.4);\else\draw(0,-.4)--(0,.4);\fi
  \draw(0,.3)--(.4,.3)--(.4,.73)node[right]{D};\draw(0,-.3)--(.4,-.3)--(.4,-.73)node[right]{S};
  \ifnum\row=3
   \draw(-.85,0)--(-.08,0);
   \ifnum\pol=0\draw[-{Triangle[length=1.9mm,width=1.4mm]}](.02,-.3)--(.38,-.3);\else\draw[-{Triangle[length=1.9mm,width=1.4mm]}](.38,-.3)--(.02,-.3);\fi
  \else
   \ifnum\pol=1\draw(-.85,0)--(-.25,0);\draw(-.165,0)circle(.085);\else\draw(-.85,0)--(-.08,0);\fi
  \fi
  \node[above]at(-.77,0){G};
 \fi
 \end{scope}}
}
\node at(1.35,4.65){Indukovani kanal};\node at(6.75,4.65){Ugrađeni kanal};
\foreach\row in{1,2,3,4}{\node at(-1.65,3.55-\row*2.2){\row};}
\end{tikzpicture}
\par\medskip D — drejn\qquad G — gejt\qquad S — sors\qquad B — osnova
\end{frame}

\slajdid{09_1-s003}
\begin{frame}[label=09_1-s003]{NMOS sa indukovanim dugačkim kanalom}\setlength{\parskip}{2pt}\setlength{\abovedisplayskip}{4pt}\setlength{\belowdisplayskip}{4pt}\setlength{\abovedisplayshortskip}{3pt}\setlength{\belowdisplayshortskip}{3pt}
% element: 09_1-s003:e01
% element: 09_1-s003:e02
% element: 09_1-s003:e03
% element: 09_1-s003:e04
% element: 09_1-s003:e05
\begin{columns}[T,onlytextwidth]\column{.56\textwidth}\centering \begin{tikzpicture}[x=.8cm,y=.65cm,font=\sitneoznake,>=Latex,line width=.7pt]
\foreach\x in{0,3}{\begin{scope}[shift={(\x,0)}]
\draw(-.6,0)--(-.1,0);\draw(-.1,-.35)--(-.1,.35);\draw(0,-.4)--(0,.4);\draw(0,.3)--(.4,.3)--(.4,.9)node[above]{D};\draw(0,-.3)--(.4,-.3)--(.4,-.9)node[below]{S};\draw[-{Triangle[length=1.6mm,width=1.2mm]}](.03,-.3)--(.34,-.3);\ifnum\x=3\node[left]at(-.9,0){G};\else\node[left]at(-.6,0){G};\fi\end{scope}}
\draw[->](-.7,.2)--(-.2,.2)node[midway,above]{$I_G$};\draw[->](.65,.85)--(.65,.35)node[midway,right]{$I_D$};\draw[->](.65,-.35)--(.65,-.85)node[midway,right]{$I_S$};
\draw[<->](3.4,.85)to[bend right=40](2.35,.05);\node at(2.5,.95){$V_{GD}$};\node at(2.25,.24){$+$};
\draw[<->](2.35,-.1)to[bend right=40](3.4,-.85);\node at(2.5,-.9){$V_{GS}$};\node at(2.2,-.22){$+$};
\draw[<->](3.65,.8)to[bend left=20](3.65,-.8);\node[right]at(3.85,0){$V_{DS}$};\node at(3.7,.95){$+$};
\end{tikzpicture}
\par\raggedright $I_{Gn}=0$, $I_{Dn}=I_{Sn}\geq0$ za $V_{GSn}\geq V_{Tn}$.\par
\textbf{Zasićenje (aktivna oblast)}\par $V_{DSn}\geq V_{GSn}-V_{Tn}$
\[I_{Dn}=\frac{k_n}{2}(V_{GSn}-V_{Tn})^2(1+\lambda_nV_{DSn})\]
\textbf{Triodna (omska, neaktivna) oblast}\par $V_{DSn}<V_{GSn}-V_{Tn}$
\[I_{Dn}=\frac{k_n}{2}(2V_{DSn}(V_{GSn}-V_{Tn})-V_{DSn}^2)\]
\column{.40\textwidth}\textbf{Parametri modela}\par
$V_{Tn}>0$, $[\lambda_n]=\mathrm{V}^{-1}$\par $W_n$, $L_n$: širina i dužina.\par $[\mu_n]=\mathrm{m^2/(Vs)}$\par $[C_{oxn}]=\mathrm{F/m^2}$
\[k_n=\mu_n C_{oxn}\frac{W_n}{L_n},\qquad [k_n]=\mathrm{A/V^2}\]
$v=\mu E$\end{columns}
\end{frame}

\slajdid{09_1-s004}
\begin{frame}[label=09_1-s004]{Modulacija dužine kanala i granica modela}\setlength{\parskip}{2pt}\setlength{\abovedisplayskip}{4pt}\setlength{\belowdisplayskip}{4pt}\setlength{\abovedisplayshortskip}{3pt}\setlength{\belowdisplayshortskip}{3pt}
% element: 09_1-s004:e01
% element: 09_1-s004:e02
% element: 09_1-s004:e03
% element: 09_1-s004:e04
% element: 09_1-s004:e05
Često zanemarujemo modulaciju: $\lambda_n=\lambda_p=0$.\par Na granici $V_{DSn}=V_{GSn}-V_{Tn}$ zasićeni izraz daje; triodni izraz nema faktor modulacije:\[I_{Dn,\mathrm{sat}}=\frac{k_n}{2}V_{DSn}^2(1+\lambda_nV_{DSn})\]

\textbf{Primeri modifikovanih modela}\par Zasićenje:\[I_{Dn}=\frac{k_n}{2}(V_{GSn}-V_{Tn})^2\bigl(1+\lambda_n(V_{DSn}-(V_{GSn}-V_T))\bigr)\]
Triodna oblast (SPICE):\[I_{Dn}=\frac{k_n}{2}\bigl(2V_{DSn}(V_{GSn}-V_{Tn})-V_{DSn}^2\bigr)(1+\lambda_nV_{DSn})\]
\medskip\alert{Koristićemo nemodifikovane izraze, uz poznavanje ograničenja modela.}
\end{frame}

\slajdid{09_1-s005}
\begin{frame}[label=09_1-s005]{Brzina nosilaca: dugačak i kratak kanal}\setlength{\parskip}{2pt}\setlength{\abovedisplayskip}{4pt}\setlength{\belowdisplayskip}{4pt}\setlength{\abovedisplayshortskip}{3pt}\setlength{\belowdisplayshortskip}{3pt}
% element: 09_1-s005:e01
% element: 09_1-s005:e02
% element: 09_1-s005:e03
% element: 09_1-s005:e04
\begin{itemize}
\item Model dugačkog kanala daje dobre rezultate za dimenzije reda $\mu\mathrm m$.
\item Pri dimenzijama reda $\mathrm{nm}$ i velikom polju nastaje zasićenje brzine.
\item Tranzistor ulazi u zasićenje pri manjem naponu $V_{DS}$.
\end{itemize}
\begin{columns}[T,onlytextwidth]\column{.44\textwidth}\centering \begin{tikzpicture}[x=.65cm,y=.65cm,font=\sitneoznake,>=Latex,line width=.7pt]\draw[->](-.15,0)--(2.5,0)node[below]{$E$};\draw[->](0,-.15)--(0,2.3)node[left]{$v$};\draw[<->](.85,0)arc[start angle=0,end angle=50,radius=.85];\node at(.35,.17){$\mu$};\draw(0,0)--(1.45,1.8);\node[align=center]at(1.35,2.4){dugačak\\kanal};\end{tikzpicture}
\hspace{5mm}\begin{tikzpicture}[x=.65cm,y=.65cm,font=\sitneoznake,>=Latex,line width=.7pt]\draw[->](-.15,0)--(2.5,0)node[below]{$E$};\draw[->](0,-.15)--(0,2.3)node[left]{$v$};\draw[<->](.85,0)arc[start angle=0,end angle=50,radius=.85];\node at(.35,.17){$\mu$};\draw(0,0)..controls(.7,.9)and(.9,1.05)..(1.1,1.1)..controls(1.5,1.2)and(1.9,1.2)..(2.3,1.2);\draw[dashed](0,0)--(1.2,1.6);\draw[dashed](0,1.2)node[left]{$v_{sat}$}--(2.3,1.2);\draw[dashed](1.1,0)node[below]{$E_C$}--(1.1,1.6);\node[align=center]at(1.5,2.4){kratak\\kanal};\end{tikzpicture}
\column{.52\textwidth}\[v_n=\mu_n\frac E{1+E/E_{Cn}}\quad(E\leq E_{Cn})\]
\[v_n=v_{nsat}\quad(E\geq E_{Cn})\]
\[v_{nsat}=\frac{\mu_nE_{Cn}}2\]
$E$: električno polje; $E_{Cn}$: kritično polje.\end{columns}
\end{frame}

\slajdid{09_1-s006}
\begin{frame}[label=09_1-s006]{Efekat kratkog kanala}\setlength{\parskip}{2pt}\setlength{\abovedisplayskip}{4pt}\setlength{\belowdisplayskip}{4pt}\setlength{\abovedisplayshortskip}{3pt}\setlength{\belowdisplayshortskip}{3pt}
% element: 09_1-s006:e01
% element: 09_1-s006:e02
% element: 09_1-s006:e03
% element: 09_1-s006:e04
\begin{columns}[c,onlytextwidth]\column{.16\textwidth}Omska, triodna\column{.82\textwidth}
\[I_{Dn}=\frac{k_n}{2}\frac1{1+\alert{\frac{V_{DSn}}{L_nE_{Cn}}}}\bigl(2V_{DSn}(V_{GSn}-V_{Tn})-V_{DSn}^2\bigr)\]
\hfill $\alert{L_n}$ veliko – svodi se na stare izraze
\end{columns}\medskip
\begin{columns}[c,onlytextwidth]\column{.16\textwidth}Zasićenje\column{.82\textwidth}
\[I_{Dn}=\frac{k_n}{2}\frac1{1+\alert{\frac{(V_{GSn}-V_{Tn})}{L_nE_{Cn}}}}(V_{GSn}-V_{Tn})^2\qquad\times(1+\lambda_nV_{DSn})\ \text{ako treba}\]
\end{columns}\medskip
\begin{columns}[c,onlytextwidth]\column{.16\textwidth}Uslov\column{.82\textwidth}
\[V_{DSn}\geq\frac{(V_{GSn}-V_{Tn})}{1+\alert{\frac{(V_{GSn}-V_{Tn})}{L_nE_{Cn}}}}\quad\textit{zasićenje}\]
\[V_{DSn}\leq\frac{(V_{GSn}-V_{Tn})}{1+\alert{\frac{(V_{GSn}-V_{Tn})}{L_nE_{Cn}}}}\quad\textit{triodna}\]
\end{columns}
\end{frame}

\slajdid{09_1-s007}
\begin{frame}[label=09_1-s007]{Parametri modela kratkog kanala}\setlength{\parskip}{2pt}\setlength{\abovedisplayskip}{4pt}\setlength{\belowdisplayskip}{4pt}\setlength{\abovedisplayshortskip}{3pt}\setlength{\belowdisplayshortskip}{3pt}
% element: 09_1-s007:e01
% element: 09_1-s007:e02
% element: 09_1-s007:e03
% element: 09_1-s007:e04
\[I_{Dn}=W_nC_{OXn}v_{nsat}\frac{(V_{GSn}-V_{Tn})^2}{(V_{GSn}-V_{Tn})+L_nE_{Cn}}\]
\[v_{nsat}=v_{psat}=v_{sat}\approx8\times10^4\frac ms.\]
\begin{columns}[T,onlytextwidth]\column{.48\textwidth}\[v_{nsat}=\frac{\mu_nE_{Cn}}2\]
\[v_{psat}=\frac{\mu_pE_{Cp}}2\]
\column{.48\textwidth}\[k_n=\mu_nC_{oxn}\frac{W_n}{L_n}\]
\[k_p=\mu_pC_{oxp}\frac{W_p}{L_p}\]
\end{columns}\medskip Ovaj zapis olakšava skraćivanje izraza u daljoj analizi.
\end{frame}

\slajdid{09_1-s008}
\begin{frame}[label=09_1-s008]{Dodatna aproksimacija brzine nosilaca}\setlength{\parskip}{2pt}\setlength{\abovedisplayskip}{4pt}\setlength{\belowdisplayskip}{4pt}\setlength{\abovedisplayshortskip}{3pt}\setlength{\belowdisplayshortskip}{3pt}
% element: 09_1-s008:e01
% element: 09_1-s008:e02
% element: 09_1-s008:e03
% element: 09_1-s008:e04
Pretpostavka: $(V_{GSn}-V_{Tn})\gg L_nE_{Cn}$.\[v_n=\mu_nE\ (E\leq E_{Cn}),\qquad v_n=v_{nsat}\ (E\geq E_{Cn})\]
\par $V_{DSnsat}=L_nE_{Cn}=L_nv_{nsat}/\mu_n$.
\begin{center}\begin{tikzpicture}[x=.65cm,y=.65cm,font=\sitneoznake,>=Latex,line width=.7pt]\draw[->](-.15,0)--(2.5,0)node[below]{$E$};\draw[->](0,-.15)--(0,2.3)node[left]{$v$};\draw[<->](.85,0)arc[start angle=0,end angle=50,radius=.85];\node at(.35,.17){$\mu$};\draw(0,0)--(1.45,1.8);\node[align=center]at(1.35,2.4){dugačak\\kanal};\end{tikzpicture}
\hspace{7mm}\begin{tikzpicture}[x=.65cm,y=.65cm,font=\sitneoznake,>=Latex,line width=.7pt]\draw[->](-.15,0)--(2.5,0)node[below]{$E$};\draw[->](0,-.15)--(0,2.3)node[left]{$v$};\draw[<->](.85,0)arc[start angle=0,end angle=50,radius=.85];\node at(.35,.17){$\mu$};\draw(0,0)..controls(.7,.9)and(.9,1.05)..(1.1,1.1)..controls(1.5,1.2)and(1.9,1.2)..(2.3,1.2);\draw[dashed](0,0)--(1.2,1.6);\draw[dashed](0,1.2)node[left]{$v_{sat}$}--(2.3,1.2);\draw[dashed](1.1,0)node[below]{$E_C$}--(1.1,1.6);\node[align=center]at(1.5,2.4){kratak\\kanal};\end{tikzpicture}
\hspace{7mm}\begin{tikzpicture}[x=.65cm,y=.65cm,font=\sitneoznake,>=Latex,line width=.7pt]\draw[->](-.15,0)--(2.5,0)node[below]{$E$};\draw[->](0,-.15)--(0,2.3)node[left]{$v$};\draw[<->](.85,0)arc[start angle=0,end angle=50,radius=.85];\node at(.35,.17){$\mu$};\draw(0,0)--(.95,1.2)--(2.3,1.2);\draw[dashed](0,1.2)node[left]{$v_{sat}$}--(.95,1.2);\draw[dashed](.95,0)node[below]{$E_C$}--(.95,1.2);\node[align=center]at(1.8,2.45){kratak kanal\\dodatna\\aproksimacija};\end{tikzpicture}
\end{center}

\[I_{Dn}=W_nC_{oxn}v_{nsat}\left(V_{GSn}-V_{Tn}-\frac{V_{DSnsat}}2\right)\]

\end{frame}

\slajdid{09_1-s009}
\begin{frame}[label=09_1-s009]{PUN i PDN mreže}\setlength{\parskip}{2pt}\setlength{\abovedisplayskip}{4pt}\setlength{\belowdisplayskip}{4pt}\setlength{\abovedisplayshortskip}{3pt}\setlength{\belowdisplayshortskip}{3pt}
% element: 09_1-s009:e01
% element: 09_1-s009:e02
% element: 09_1-s009:e03
% element: 09_1-s009:e04
\begin{columns}[T,onlytextwidth]\column{.34\textwidth}\centering \begin{tikzpicture}[x=.8cm,y=.7cm,font=\sitneoznake,>=Latex,line width=.7pt]
\draw(0,0)rectangle(2,3.6);\draw(.4,0)--(.4,3.6);\draw(.4,1.65)--(2,1.65);\draw(.4,1.95)--(2,1.95);\node[align=center]at(1.2,2.7){Pull up\\network};\node[align=center]at(1.2,.8){Pull down\\network};\draw(1.05,3.6)--(1.05,4);\draw(.8,4)--(1.3,4)node[above left]{$V_{DD}$};\draw(1.05,0)--(1.05,-.4)node[ground]{};
\draw[->](1.85,1.8)--(2.65,1.8)node[right]{Izlaz};\draw(-.75,1.98)--(-.3,1.98);\draw(-.75,1.62)--(-.3,1.62);\draw[->](-.8,1.8)--(0,1.8);\node[left]at(-.8,1.8){Ulazi};
\end{tikzpicture}
\column{.62\textwidth}\textbf{PUN — pull up network}\par Deo kola koji obezbeđuje logičku jedinicu.
\medskip\par\textbf{PDN — pull down network}\par Deo kola koji obezbeđuje logičku nulu.
\medskip\par\textbf{Opterećenje na izlazu}\par Pasivno: PUN je otpornik.\par Aktivno: PUN čine aktivni elementi, tranzistori.\end{columns}\medskip
\begin{itemize}
\item Cilj: komplementarne mreže — dok jedna radi, druga je isključena.
\item \alert{Trostatička kola}: obe mreže mogu biti isključene.
\item \alert{Otvoreni drejn}: kolo ima samo PDN.
\end{itemize}
\end{frame}

\slajdid{09_1-s010}
\begin{frame}[label=09_1-s010]{PDN — moguće realizacije}\setlength{\parskip}{2pt}\setlength{\abovedisplayskip}{4pt}\setlength{\belowdisplayskip}{4pt}\setlength{\abovedisplayshortskip}{3pt}\setlength{\belowdisplayshortskip}{3pt}
% element: 09_1-s010:e01
% element: 09_1-s010:e02
% element: 09_1-s010:e03
% element: 09_1-s010:e04
\begin{columns}[T,onlytextwidth]\column{.25\textwidth}\centering \begin{tikzpicture}[x=.58cm,y=.6cm,font=\sitneoznake,>=Latex,line width=.7pt]
\draw(0,0)rectangle(2.9,3.6);\draw(.4,0)--(.4,3.6);\draw(.4,1.65)--(2.9,1.65);\draw(.4,1.95)--(2.9,1.95);\node[align=center]at(1.65,2.7){Pull up\\network};\node[align=center]at(1.65,.8){Pull down\\network};\draw(1.05,3.6)--(1.05,4);\draw(.8,4)--(1.3,4)node[above left]{$V_{DD}$};\draw(1.05,0)--(1.05,-.4)node[ground]{};
\draw[->](2.75,1.8)--(3.2,1.8)node[right]{Izl.};\draw(-.75,1.98)--(-.3,1.98);\draw(-.75,1.62)--(-.3,1.62);\draw[->](-.6,1.8)--(0,1.8);\node[left]at(-.6,1.8){Ul.};
\end{tikzpicture}

\column{.71\textwidth}\textbf{Paralelno — NILI}\par $Y=\overline{A+B}$\par\smallskip\begin{columns}[c,onlytextwidth]\column{.31\textwidth}\centering \begin{tikzpicture}[x=.58cm,y=.36cm,font=\sitneoznake,>=Latex,line width=.7pt]\begin{scope}[shift={(0,0)}]\draw(0,.55)--(0,.25);\draw(0,-.25)--(0,-.55);\draw(0,-.25)--(-.12,.2);\draw(-.7,0)--(-.15,0);\node[left]at(-.7,0){$A$};\end{scope}\begin{scope}[shift={(1.45,0)}]\draw(0,.55)--(0,.25);\draw(0,-.25)--(0,-.55);\draw(0,-.25)--(-.12,.2);\draw(-.7,0)--(-.15,0);\node[left]at(-.7,0){$B$};\end{scope}\draw(0,.55)--(0,.9)--(1.45,.9)--(1.45,.55);\draw(0,-.55)--(0,-.85)--(1.45,-.85)--(1.45,-.55);\draw(0.725,0.9)--(0.725,1.12)--(0.815,1.175)--(0.635,1.285)--(0.815,1.395)--(0.635,1.505)--(0.815,1.615)--(0.635,1.725)--(0.725,1.78)--(0.725,2);\draw(0.725,2)--(0.725,2.2);\draw(.5,2.2)--(.95,2.2)node[above left]{$V_{DD}$};\draw(0.725,0.9)--(1.9,0.9)node[right]{$Y$};\draw(0.725,-0.85)--(0.725,-1.35)node[ground]{};\foreach\p in {(0.725,.9),(0.725,-.85)}{\fill\p circle(1pt);}
\end{tikzpicture}
\column{.31\textwidth}\centering \begin{tikzpicture}[x=.58cm,y=.36cm,font=\sitneoznake,>=Latex,line width=.7pt]\begin{scope}[shift={(0,0)}]\draw(0,.55)--(0,.23)--(-.22,.23);\draw(0,-.55)--(0,-.23)--(-.22,-.23);\draw(-.22,-.32)--(-.22,.32);\draw(-.28,-.28)--(-.28,.28);\draw(-.7,0)--(-.28,0);\draw[-{Triangle[length=1.6mm,width=1.3mm]}](-.21,-.23)--(0,-.23);\node[left]at(-.7,0){$A$};\end{scope}\begin{scope}[shift={(1.45,0)}]\draw(0,.55)--(0,.23)--(-.22,.23);\draw(0,-.55)--(0,-.23)--(-.22,-.23);\draw(-.22,-.32)--(-.22,.32);\draw(-.28,-.28)--(-.28,.28);\draw(-.7,0)--(-.28,0);\draw[-{Triangle[length=1.6mm,width=1.3mm]}](-.21,-.23)--(0,-.23);\node[left]at(-.7,0){$B$};\end{scope}\draw(0,.55)--(0,.9)--(1.45,.9)--(1.45,.55);\draw(0,-.55)--(0,-.85)--(1.45,-.85)--(1.45,-.55);\draw(0.725,0.9)--(0.725,1.12)--(0.815,1.175)--(0.635,1.285)--(0.815,1.395)--(0.635,1.505)--(0.815,1.615)--(0.635,1.725)--(0.725,1.78)--(0.725,2);\draw(0.725,2)--(0.725,2.2);\draw(.5,2.2)--(.95,2.2)node[above left]{$V_{DD}$};\draw(0.725,0.9)--(1.9,0.9)node[right]{$Y$};\draw(0.725,-0.85)--(0.725,-1.35)node[ground]{};\foreach\p in {(0.725,.9),(0.725,-.85)}{\fill\p circle(1pt);}
\end{tikzpicture}
\column{.31\textwidth}\centering \begin{tikzpicture}[x=.58cm,y=.36cm,font=\sitneoznake,>=Latex,line width=.7pt]\begin{scope}[shift={(0,0)}]\draw(0,.55)--(0,.23)--(-.22,.23);\draw(0,-.55)--(0,-.23)--(-.22,-.23);\draw(-.22,-.32)--(-.22,.32);\draw(-.28,-.28)--(-.28,.28);\draw(-.7,0)--(-.28,0);\node[left]at(-.7,0){$A$};\end{scope}\begin{scope}[shift={(1.45,0)}]\draw(0,.55)--(0,.23)--(-.22,.23);\draw(0,-.55)--(0,-.23)--(-.22,-.23);\draw(-.22,-.32)--(-.22,.32);\draw(-.28,-.28)--(-.28,.28);\draw(-.7,0)--(-.28,0);\node[left]at(-.7,0){$B$};\end{scope}\draw(0,.55)--(0,.9)--(1.45,.9)--(1.45,.55);\draw(0,-.55)--(0,-.85)--(1.45,-.85)--(1.45,-.55);\draw(0.725,0.9)--(0.725,1.12)--(0.815,1.175)--(0.635,1.285)--(0.815,1.395)--(0.635,1.505)--(0.815,1.615)--(0.635,1.725)--(0.725,1.78)--(0.725,2);\draw(0.725,2)--(0.725,2.2);\draw(.5,2.2)--(.95,2.2)node[above left]{$V_{DD}$};\draw(0.725,0.9)--(1.9,0.9)node[right]{$Y$};\draw(0.725,-0.85)--(0.725,-1.35)node[ground]{};\foreach\p in {(0.725,.9),(0.725,-.85)}{\fill\p circle(1pt);}
\end{tikzpicture}
\end{columns}
\smallskip\textbf{Redno — NI}\par $Y=\overline{AB}$\par\smallskip\begin{columns}[c,onlytextwidth]\column{.31\textwidth}\centering \begin{tikzpicture}[x=.58cm,y=.36cm,font=\sitneoznake,>=Latex,line width=.7pt]\begin{scope}[shift={(0.725,0.55)}]\draw(0,.55)--(0,.25);\draw(0,-.25)--(0,-.55);\draw(0,-.25)--(-.12,.2);\draw(-.7,0)--(-.15,0);\node[left]at(-.7,0){$A$};\end{scope}\begin{scope}[shift={(0.725,-0.55)}]\draw(0,.55)--(0,.25);\draw(0,-.25)--(0,-.55);\draw(0,-.25)--(-.12,.2);\draw(-.7,0)--(-.15,0);\node[left]at(-.7,0){$B$};\end{scope}\draw(0.725,1.1)--(0.725,1.28)--(0.815,1.325)--(0.635,1.415)--(0.815,1.505)--(0.635,1.595)--(0.815,1.685)--(0.635,1.775)--(0.725,1.82)--(0.725,2);\draw(0.725,2)--(0.725,2.2);\draw(.5,2.2)--(.95,2.2)node[above left]{$V_{DD}$};\draw(0.725,1.1)--(1.9,1.1)node[right]{$Y$};\draw(0.725,-1.1)--(0.725,-1.35)node[ground]{};\end{tikzpicture}
\column{.31\textwidth}\centering \begin{tikzpicture}[x=.58cm,y=.36cm,font=\sitneoznake,>=Latex,line width=.7pt]\begin{scope}[shift={(0.725,0.55)}]\draw(0,.55)--(0,.23)--(-.22,.23);\draw(0,-.55)--(0,-.23)--(-.22,-.23);\draw(-.22,-.32)--(-.22,.32);\draw(-.28,-.28)--(-.28,.28);\draw(-.7,0)--(-.28,0);\draw[-{Triangle[length=1.6mm,width=1.3mm]}](-.21,-.23)--(0,-.23);\node[left]at(-.7,0){$A$};\end{scope}\begin{scope}[shift={(0.725,-0.55)}]\draw(0,.55)--(0,.23)--(-.22,.23);\draw(0,-.55)--(0,-.23)--(-.22,-.23);\draw(-.22,-.32)--(-.22,.32);\draw(-.28,-.28)--(-.28,.28);\draw(-.7,0)--(-.28,0);\draw[-{Triangle[length=1.6mm,width=1.3mm]}](-.21,-.23)--(0,-.23);\node[left]at(-.7,0){$B$};\end{scope}\draw(0.725,1.1)--(0.725,1.28)--(0.815,1.325)--(0.635,1.415)--(0.815,1.505)--(0.635,1.595)--(0.815,1.685)--(0.635,1.775)--(0.725,1.82)--(0.725,2);\draw(0.725,2)--(0.725,2.2);\draw(.5,2.2)--(.95,2.2)node[above left]{$V_{DD}$};\draw(0.725,1.1)--(1.9,1.1)node[right]{$Y$};\draw(0.725,-1.1)--(0.725,-1.35)node[ground]{};\end{tikzpicture}
\column{.31\textwidth}\centering \begin{tikzpicture}[x=.58cm,y=.36cm,font=\sitneoznake,>=Latex,line width=.7pt]\begin{scope}[shift={(0.725,0.55)}]\draw(0,.55)--(0,.23)--(-.22,.23);\draw(0,-.55)--(0,-.23)--(-.22,-.23);\draw(-.22,-.32)--(-.22,.32);\draw(-.28,-.28)--(-.28,.28);\draw(-.7,0)--(-.28,0);\node[left]at(-.7,0){$A$};\end{scope}\begin{scope}[shift={(0.725,-0.55)}]\draw(0,.55)--(0,.23)--(-.22,.23);\draw(0,-.55)--(0,-.23)--(-.22,-.23);\draw(-.22,-.32)--(-.22,.32);\draw(-.28,-.28)--(-.28,.28);\draw(-.7,0)--(-.28,0);\node[left]at(-.7,0){$B$};\end{scope}\draw(0.725,1.1)--(0.725,1.28)--(0.815,1.325)--(0.635,1.415)--(0.815,1.505)--(0.635,1.595)--(0.815,1.685)--(0.635,1.775)--(0.725,1.82)--(0.725,2);\draw(0.725,2)--(0.725,2.2);\draw(.5,2.2)--(.95,2.2)node[above left]{$V_{DD}$};\draw(0.725,1.1)--(1.9,1.1)node[right]{$Y$};\draw(0.725,-1.1)--(0.725,-1.35)node[ground]{};\end{tikzpicture}
\end{columns}
\end{columns}
\end{frame}

\slajdid{09_1-s011}
\begin{frame}[label=09_1-s011]{PUN — moguće realizacije}\setlength{\parskip}{2pt}\setlength{\abovedisplayskip}{4pt}\setlength{\belowdisplayskip}{4pt}\setlength{\abovedisplayshortskip}{3pt}\setlength{\belowdisplayshortskip}{3pt}
% element: 09_1-s011:e01
% element: 09_1-s011:e02
% element: 09_1-s011:e03
% element: 09_1-s011:e04
\begin{columns}[T,onlytextwidth]\column{.25\textwidth}\centering \begin{tikzpicture}[x=.58cm,y=.6cm,font=\sitneoznake,>=Latex,line width=.7pt]
\draw(0,0)rectangle(2.9,3.6);\draw(.4,0)--(.4,3.6);\draw(.4,1.65)--(2.9,1.65);\draw(.4,1.95)--(2.9,1.95);\node[align=center]at(1.65,2.7){Pull up\\network};\node[align=center]at(1.65,.8){Pull down\\network};\draw(1.05,3.6)--(1.05,4);\draw(.8,4)--(1.3,4)node[above left]{$V_{DD}$};\draw(1.05,0)--(1.05,-.4)node[ground]{};
\draw[->](2.75,1.8)--(3.2,1.8)node[right]{Izl.};\draw(-.75,1.98)--(-.3,1.98);\draw(-.75,1.62)--(-.3,1.62);\draw[->](-.6,1.8)--(0,1.8);\node[left]at(-.6,1.8){Ul.};
\end{tikzpicture}

\column{.71\textwidth}\textbf{Paralelno — ILI sa aktivnim nulama; dualno NI}\par $Y=\overline A+\overline B=\overline{AB}$\par\smallskip\begin{columns}[c,onlytextwidth]\column{.31\textwidth}\centering \begin{tikzpicture}[x=.58cm,y=.36cm,font=\sitneoznake,>=Latex,line width=.7pt]\begin{scope}[shift={(0,0)}]\draw(0,.55)--(0,.25);\draw(0,-.25)--(0,-.55);\draw(0,-.25)--(-.12,.2);\draw(-.7,0)--(-.15,0);\draw[fill=white](-.23,0)circle(.065);\node[left]at(-.7,0){$A$};\end{scope}\begin{scope}[shift={(1.45,0)}]\draw(0,.55)--(0,.25);\draw(0,-.25)--(0,-.55);\draw(0,-.25)--(-.12,.2);\draw(-.7,0)--(-.15,0);\draw[fill=white](-.23,0)circle(.065);\node[left]at(-.7,0){$B$};\end{scope}\draw(0,.55)--(0,.9)--(1.45,.9)--(1.45,.55);\draw(0,-.55)--(0,-.85)--(1.45,-.85)--(1.45,-.55);\draw(0.725,0.9)--(0.725,1.55);\draw(.5,1.55)--(.95,1.55)node[above left]{$V_{DD}$};\draw(0.725,-0.85)--(1.9,-0.85)node[right]{$Y$};\draw(0.725,-0.85)--(0.725,-1.09)--(0.815,-1.15)--(0.635,-1.27)--(0.815,-1.39)--(0.635,-1.51)--(0.815,-1.63)--(0.635,-1.75)--(0.725,-1.81)--(0.725,-2.05);\draw(0.725,-2.05)--(0.725,-2.2)node[ground]{};\foreach\p in {(0.725,.9),(0.725,-.85)}{\fill\p circle(1pt);}
\end{tikzpicture}
\column{.31\textwidth}\centering \begin{tikzpicture}[x=.58cm,y=.36cm,font=\sitneoznake,>=Latex,line width=.7pt]\begin{scope}[shift={(0,0)}]\draw(0,.55)--(0,.23)--(-.22,.23);\draw(0,-.55)--(0,-.23)--(-.22,-.23);\draw(-.22,-.32)--(-.22,.32);\draw(-.28,-.28)--(-.28,.28);\draw(-.7,0)--(-.28,0);\draw[-{Triangle[length=1.6mm,width=1.3mm]}](0,.23)--(-.21,.23);\node[left]at(-.7,0){$A$};\end{scope}\begin{scope}[shift={(1.45,0)}]\draw(0,.55)--(0,.23)--(-.22,.23);\draw(0,-.55)--(0,-.23)--(-.22,-.23);\draw(-.22,-.32)--(-.22,.32);\draw(-.28,-.28)--(-.28,.28);\draw(-.7,0)--(-.28,0);\draw[-{Triangle[length=1.6mm,width=1.3mm]}](0,.23)--(-.21,.23);\node[left]at(-.7,0){$B$};\end{scope}\draw(0,.55)--(0,.9)--(1.45,.9)--(1.45,.55);\draw(0,-.55)--(0,-.85)--(1.45,-.85)--(1.45,-.55);\draw(0.725,0.9)--(0.725,1.55);\draw(.5,1.55)--(.95,1.55)node[above left]{$V_{DD}$};\draw(0.725,-0.85)--(1.9,-0.85)node[right]{$Y$};\draw(0.725,-0.85)--(0.725,-1.09)--(0.815,-1.15)--(0.635,-1.27)--(0.815,-1.39)--(0.635,-1.51)--(0.815,-1.63)--(0.635,-1.75)--(0.725,-1.81)--(0.725,-2.05);\draw(0.725,-2.05)--(0.725,-2.2)node[ground]{};\foreach\p in {(0.725,.9),(0.725,-.85)}{\fill\p circle(1pt);}
\end{tikzpicture}
\column{.31\textwidth}\centering \begin{tikzpicture}[x=.58cm,y=.36cm,font=\sitneoznake,>=Latex,line width=.7pt]\begin{scope}[shift={(0,0)}]\draw(0,.55)--(0,.23)--(-.22,.23);\draw(0,-.55)--(0,-.23)--(-.22,-.23);\draw(-.22,-.32)--(-.22,.32);\draw(-.28,-.28)--(-.28,.28);\draw(-.7,0)--(-.43,0);\draw(-.355,0)circle(.075);\node[left]at(-.7,0){$A$};\end{scope}\begin{scope}[shift={(1.45,0)}]\draw(0,.55)--(0,.23)--(-.22,.23);\draw(0,-.55)--(0,-.23)--(-.22,-.23);\draw(-.22,-.32)--(-.22,.32);\draw(-.28,-.28)--(-.28,.28);\draw(-.7,0)--(-.43,0);\draw(-.355,0)circle(.075);\node[left]at(-.7,0){$B$};\end{scope}\draw(0,.55)--(0,.9)--(1.45,.9)--(1.45,.55);\draw(0,-.55)--(0,-.85)--(1.45,-.85)--(1.45,-.55);\draw(0.725,0.9)--(0.725,1.55);\draw(.5,1.55)--(.95,1.55)node[above left]{$V_{DD}$};\draw(0.725,-0.85)--(1.9,-0.85)node[right]{$Y$};\draw(0.725,-0.85)--(0.725,-1.09)--(0.815,-1.15)--(0.635,-1.27)--(0.815,-1.39)--(0.635,-1.51)--(0.815,-1.63)--(0.635,-1.75)--(0.725,-1.81)--(0.725,-2.05);\draw(0.725,-2.05)--(0.725,-2.2)node[ground]{};\foreach\p in {(0.725,.9),(0.725,-.85)}{\fill\p circle(1pt);}
\end{tikzpicture}
\end{columns}
\textbf{Redno — I sa aktivnim nulama; dualno NILI}\par $Y=\overline A\,\overline B=\overline{A+B}$\par\smallskip\begin{columns}[c,onlytextwidth]\column{.31\textwidth}\centering \begin{tikzpicture}[x=.58cm,y=.36cm,font=\sitneoznake,>=Latex,line width=.7pt]\begin{scope}[shift={(0.725,0.55)}]\draw(0,.55)--(0,.25);\draw(0,-.25)--(0,-.55);\draw(0,-.25)--(-.12,.2);\draw(-.7,0)--(-.15,0);\draw[fill=white](-.23,0)circle(.065);\node[left]at(-.7,0){$A$};\end{scope}\begin{scope}[shift={(0.725,-0.55)}]\draw(0,.55)--(0,.25);\draw(0,-.25)--(0,-.55);\draw(0,-.25)--(-.12,.2);\draw(-.7,0)--(-.15,0);\draw[fill=white](-.23,0)circle(.065);\node[left]at(-.7,0){$B$};\end{scope}\draw(0.725,1.1)--(0.725,1.55);\draw(.5,1.55)--(.95,1.55)node[above left]{$V_{DD}$};\draw(0.725,-1.1)--(1.9,-1.1)node[right]{$Y$};\draw(0.725,-1.1)--(0.725,-1.29)--(0.815,-1.3375)--(0.635,-1.4325)--(0.815,-1.5275)--(0.635,-1.6225)--(0.815,-1.7175)--(0.635,-1.8125)--(0.725,-1.86)--(0.725,-2.05);\draw(0.725,-2.05)--(0.725,-2.2)node[ground]{};\end{tikzpicture}
\column{.31\textwidth}\centering \begin{tikzpicture}[x=.58cm,y=.36cm,font=\sitneoznake,>=Latex,line width=.7pt]\begin{scope}[shift={(0.725,0.55)}]\draw(0,.55)--(0,.23)--(-.22,.23);\draw(0,-.55)--(0,-.23)--(-.22,-.23);\draw(-.22,-.32)--(-.22,.32);\draw(-.28,-.28)--(-.28,.28);\draw(-.7,0)--(-.28,0);\draw[-{Triangle[length=1.6mm,width=1.3mm]}](0,.23)--(-.21,.23);\node[left]at(-.7,0){$A$};\end{scope}\begin{scope}[shift={(0.725,-0.55)}]\draw(0,.55)--(0,.23)--(-.22,.23);\draw(0,-.55)--(0,-.23)--(-.22,-.23);\draw(-.22,-.32)--(-.22,.32);\draw(-.28,-.28)--(-.28,.28);\draw(-.7,0)--(-.28,0);\draw[-{Triangle[length=1.6mm,width=1.3mm]}](0,.23)--(-.21,.23);\node[left]at(-.7,0){$B$};\end{scope}\draw(0.725,1.1)--(0.725,1.55);\draw(.5,1.55)--(.95,1.55)node[above left]{$V_{DD}$};\draw(0.725,-1.1)--(1.9,-1.1)node[right]{$Y$};\draw(0.725,-1.1)--(0.725,-1.29)--(0.815,-1.3375)--(0.635,-1.4325)--(0.815,-1.5275)--(0.635,-1.6225)--(0.815,-1.7175)--(0.635,-1.8125)--(0.725,-1.86)--(0.725,-2.05);\draw(0.725,-2.05)--(0.725,-2.2)node[ground]{};\end{tikzpicture}
\column{.31\textwidth}\centering \begin{tikzpicture}[x=.58cm,y=.36cm,font=\sitneoznake,>=Latex,line width=.7pt]\begin{scope}[shift={(0.725,0.55)}]\draw(0,.55)--(0,.23)--(-.22,.23);\draw(0,-.55)--(0,-.23)--(-.22,-.23);\draw(-.22,-.32)--(-.22,.32);\draw(-.28,-.28)--(-.28,.28);\draw(-.7,0)--(-.43,0);\draw(-.355,0)circle(.075);\node[left]at(-.7,0){$A$};\end{scope}\begin{scope}[shift={(0.725,-0.55)}]\draw(0,.55)--(0,.23)--(-.22,.23);\draw(0,-.55)--(0,-.23)--(-.22,-.23);\draw(-.22,-.32)--(-.22,.32);\draw(-.28,-.28)--(-.28,.28);\draw(-.7,0)--(-.43,0);\draw(-.355,0)circle(.075);\node[left]at(-.7,0){$B$};\end{scope}\draw(0.725,1.1)--(0.725,1.55);\draw(.5,1.55)--(.95,1.55)node[above left]{$V_{DD}$};\draw(0.725,-1.1)--(1.9,-1.1)node[right]{$Y$};\draw(0.725,-1.1)--(0.725,-1.29)--(0.815,-1.3375)--(0.635,-1.4325)--(0.815,-1.5275)--(0.635,-1.6225)--(0.815,-1.7175)--(0.635,-1.8125)--(0.725,-1.86)--(0.725,-2.05);\draw(0.725,-2.05)--(0.725,-2.2)node[ground]{};\end{tikzpicture}
\end{columns}
\end{columns}
\end{frame}

\slajdid{09_1-s012}
\begin{frame}[label=09_1-s012]{Invertor sa pasivnim opterećenjem}\setlength{\parskip}{2pt}\setlength{\abovedisplayskip}{4pt}\setlength{\belowdisplayskip}{4pt}\setlength{\abovedisplayshortskip}{3pt}\setlength{\belowdisplayshortskip}{3pt}
% element: 09_1-s012:e01
% element: 09_1-s012:e02
% element: 09_1-s012:e03
% element: 09_1-s012:e04
\begin{center}{\RazaviSetup
\begin{circuitikz}[labeltext,american resistors,line width=.7pt]

\node[rz nmos] (d) at (0,0) {};
\draw[wire] (d.D)--(0,.65) to[american resistor,l=$R_L$] (0,1.65);
\coordinate (supply) at (0,1.65);\RazaviSupply{supply}{V_{DD}}
\draw[wire] (d.S)--(0,-.7);\coordinate(gnd)at(0,-.7);\RazaviGround{gnd}
\draw[wire] (d.G)--++(-.35,0) node[rz terminal]{} node[left=2pt]{$V_I$};
\draw[wire](0,.65) node[junction]{} --(.55,.65)node[rz terminal]{}node[right=2pt]{$V_O$};
\node[right=2pt]at(d.center){$T_D$};
\end{circuitikz}}
\hspace{15mm}% Kvalitativna kriva iz originala: nema uvedene numeričke skale.
\begin{tikzpicture}[DEoznaka,baseline=(current bounding box.center)]
\begin{axis}[width=25mm,height=20mm,scale only axis,
 axis lines=middle,xmin=-.08,xmax=1.22,ymin=-.06,ymax=1.2,
 ticks=none,clip=false,xlabel=\(V_I\),ylabel=\(V_O\),
 every axis x label/.style={at={(axis description cs:1.20,.03)},anchor=west},
 every axis y label/.style={at={(axis description cs:0,1)},anchor=south},
 axis line style={DEvod,-{Latex[length=1.5mm]}}]
\addplot[DEvod,smooth] coordinates {(0,1)(.13,1)(.24,.98)(.35,.85)(.43,.55)(.52,.22)(.60,.12)(.85,.085)(1,.06)};
\draw[dashed] (axis cs:.13,0)--(axis cs:.13,1.1);
\draw[dashed] (axis cs:.59,0)--(axis cs:.59,1.1);
\draw[dashed] (axis cs:1,0)--(axis cs:1,1.1);
\node[above] at (axis cs:.075,1.02) {1.};
\node[above] at (axis cs:.34,1.02) {2.};
\node[above] at (axis cs:.80,1.02) {3.};
\node[left] at (axis cs:0,1) {\(V_{DD}\)};
\node[below] at (axis cs:1,0) {\(V_{DD}\)};
\end{axis}
\end{tikzpicture}
\end{center}\textbf{Prva oblast:} $V_I<V_{Tn}$; $T_D$ ne vodi.\par $V_I=V_{GS,D}<V_{Tn}$.
\[
\begin{aligned}
V_O=V_{DD}-R_LI_{RL}=V_{DD}-R_LI_{Dn,D}=V_{DD}
\end{aligned}
\]

\end{frame}

\slajdid{09_1-s013}
\begin{frame}[label=09_1-s013]{Implicitna karakteristika prenosa}\setlength{\parskip}{2pt}\setlength{\abovedisplayskip}{4pt}\setlength{\belowdisplayskip}{4pt}\setlength{\abovedisplayshortskip}{3pt}\setlength{\belowdisplayshortskip}{3pt}
% element: 09_1-s013:e01
% element: 09_1-s013:e02
% element: 09_1-s013:e03
% element: 09_1-s013:e04
Eksplicitni zapis $V_O=f(V_I)$ nije uvek potreban.\par\textbf{Neopterećeno kolo:}\[\alert{I_{PUN}=I_{PDN}}\]
\textbf{Druga oblast — tranzistor u zasićenju}\[\alert{I_{PUN}=\frac{V_{DD}-V_O}{R_L}=I_{Dn,D}=\frac{k_n}{2}\frac1{1+\frac{(V_{GS,D}-V_{Tn})}{L_nE_{Cn}}}(V_{GS,D}-V_{Tn})^2(1+\lambda_nV_{DS,D})}\]
\medskip Za statičku analizu: $\lambda_n\approx0$, $V_{GS,D}=V_I$.\[\frac{V_{DD}-V_O}{R_L}=\frac{k_n}{2}\frac1{1+\frac{(V_I-V_{Tn})}{L_nE_{Cn}}}(V_I-V_{Tn})^2\]

\end{frame}

\slajdid{09_1-s014}
\begin{frame}[label=09_1-s014]{Granica zasićenja pogonskog tranzistora}\setlength{\parskip}{2pt}\setlength{\abovedisplayskip}{4pt}\setlength{\belowdisplayskip}{4pt}\setlength{\abovedisplayshortskip}{3pt}\setlength{\belowdisplayshortskip}{3pt}
% element: 09_1-s014:e01
% element: 09_1-s014:e02
% element: 09_1-s014:e03
% element: 09_1-s014:e04

Zavisnost liči na obrnutu, iskrivljenu parabolu; maksimum je u $V_I=V_{Tn}$.\par\medskip\textbf{Izlazak iz zasićenja}\[V_O=\frac{(V_I-V_{Tn})L_nE_{Cn}}{L_nE_{Cn}+(V_I-V_{Tn})}=\frac{(V_I-V_{Tn})}{1+\frac{(V_I-V_{Tn})}{L_nE_{Cn}}}\]
\medskip Bez efekta kratkog kanala:\[V_I=V_{Tn}+\sqrt{2\frac{V_{DD}}{k_nR_L}}\]
Uslov za pojavu treće oblasti:\[\left(V_{Tn}+\sqrt{2\frac{V_{DD}}{k_nR_L}}\right)<V_{DD}\]

\end{frame}

\slajdid{09_1-s015}
\begin{frame}[label=09_1-s015]{Oblasti rada i pojačanje}\setlength{\parskip}{2pt}\setlength{\abovedisplayskip}{4pt}\setlength{\belowdisplayskip}{4pt}\setlength{\abovedisplayshortskip}{3pt}\setlength{\belowdisplayshortskip}{3pt}
% element: 09_1-s015:e01
% element: 09_1-s015:e02
% element: 09_1-s015:e03
% element: 09_1-s015:e04
\begin{center}\begin{tikzpicture}[x=.68cm,y=.59cm,font=\sitneoznake,>=Latex,remember picture,line width=.7pt]\draw[->](-.15,0)--(3.3,0)node[right,yshift=4pt]{$V_I$};\draw[->](0,-.1)--(0,2.7)node[above]{$V_O$};\node[left]at(0,2.35){$V_{DD}$};\node[below]at(2.8,0){$V_{DD}$};\draw(0,2.35)--(.35,2.35)..controls(.65,2.35)and(.7,2.2)..(.85,1.75)..controls(1.05,1.2)and(1.25,.3)..(1.55,.25)--(2.8,.15);\draw[dashed](2.8,-.1)--(2.8,2.65);\draw[dashed](.3,-.1)--(.3,2.65);\draw[dashed,DEcrvena](1.5,-.1)--(1.5,2.65);\node at(.22,2.6){1.};\node at(.87,2.6){2.};\node at(2.1,2.6){3.};\draw[->](2.75,2.3)--(1.55,.32);\node[above]at(2.5,2.5){$V_O=V_I-V_{Tn}$};\end{tikzpicture}
\hspace{10mm}\begin{tikzpicture}[x=.68cm,y=.59cm,font=\sitneoznake,>=Latex,remember picture,line width=.7pt]\draw[->](-.15,0)--(3.3,0)node[right,yshift=4pt]{$V_I$};\draw[->](0,-.1)--(0,2.7)node[above]{$V_O$};\node[left]at(0,2.35){$V_{DD}$};\node[below]at(2.8,0){$V_{DD}$};\draw(0,2.35)--(.35,2.35)..controls(.65,2.35)and(.7,2.2)..(.85,1.75)..controls(1.05,1.2)and(1.25,.3)..(1.55,.25)--(2.8,.15);\draw[dashed](2.8,-.1)--(2.8,2.65);\draw[dashed,DEcrvena](.58,-.1)--(.58,2.65);\draw[dashed,DEcrvena](1.55,-.1)--(1.55,2.65);\coordinate(s15high)at(.58,2.31);\coordinate(s15low)at(1.55,.25);\node at(.3,2.6){1.};\node at(1.05,2.6){2.};\node at(2.15,2.6){3.};\draw[->](2.4,1.25)--(1.55,.32);\node[above]at(2.35,1.3){$a=-1$};\end{tikzpicture}
\end{center}1. Zakočen / zasićen.\quad 2. Zasićen / triodni.\quad 3. Triodni.
\par\[\frac{V_{DD}-V_O}{R_L}=\frac{k_n}{2}\frac1{1+\frac{V_O}{L_nE_{Cn}}}\bigl(2V_O(V_I-V_{Tn})-V_O^2\bigr)\]
\[\frac{V_{DD}-V_O}{R_L}=\frac{k_n}{2}\frac{(V_I-V_{Tn})^2}{1+\frac{V_I-V_{Tn}}{L_nE_{Cn}}}\]

\end{frame}

\slajdid{09_1-s016}
\begin{frame}[label=09_1-s016]{Izlazni naponski nivoi}\setlength{\parskip}{2pt}\setlength{\abovedisplayskip}{4pt}\setlength{\belowdisplayskip}{4pt}\setlength{\abovedisplayshortskip}{3pt}\setlength{\belowdisplayshortskip}{3pt}
% element: 09_1-s016:e01
% element: 09_1-s016:e02
% element: 09_1-s016:e03
\textbf{Logička jedinica}\[V_{OH}=V_{DD}\]
\textbf{Logička nula: treća oblast, $V_I=V_{OH}=V_{DD}$}\[\frac{V_{DD}-V_{OL}}{R_L}=\frac{k_n}{2}\frac1{1+\frac{V_{OL}}{L_nE_{Cn}}}\bigl(2V_{OL}(V_{DD}-V_{Tn})-V_{OL}^2\bigr)\]
\begin{itemize}
\item Jednačina se može rešiti kao kvadratna.
\item \alert{Koristićemo zanemarivanja}: kvalitativna slika uz kvantitativnu podršku.
\item \alert{Postoje situacije kada zanemarivanja nisu dozvoljena.}
\end{itemize}
\end{frame}

\slajdid{09_1-s017}
\begin{frame}[label=09_1-s017]{Aproksimacija napona logičke nule}\setlength{\parskip}{2pt}\setlength{\abovedisplayskip}{4pt}\setlength{\belowdisplayskip}{4pt}\setlength{\abovedisplayshortskip}{3pt}\setlength{\belowdisplayshortskip}{3pt}
% element: 09_1-s017:e01
% element: 09_1-s017:e02
% element: 09_1-s017:e03
% element: 09_1-s017:e04
\textbf{Pretpostavke}\[V_{OL}\approx0,\qquad V_{OL}\ll L_nE_{Cn},\qquad V_{DD}-V_{Tn}\gg V_{OL}\]
Zanemarujemo korekciju kratkog kanala i kvadratni član:\[\frac{V_{DD}-V_{OL}}{R_L}\approx k_nV_{OL}(V_{DD}-V_{Tn})\]
\[\frac{V_{DD}}{R_L}\approx V_{OL}\left(k_n(V_{DD}-V_{Tn})+\frac1{R_L}\right)\]
\medskip\alert{\[V_{OL}\approx\frac{V_{DD}}{k_nR_L(V_{DD}-V_{Tn})+1}\]
}
\end{frame}

\slajdid{09_1-s018}
\begin{frame}[label=09_1-s018]{Određivanje ulaznih pragova}\setlength{\parskip}{2pt}\setlength{\abovedisplayskip}{4pt}\setlength{\belowdisplayskip}{4pt}\setlength{\abovedisplayshortskip}{3pt}\setlength{\belowdisplayshortskip}{3pt}
% element: 09_1-s018:e01
% element: 09_1-s018:e02
% element: 09_1-s018:e03
% element: 09_1-s018:e04
Pragove nalazimo iz definicije: pojačanje je po apsolutnoj vrednosti 1.\[\left|\frac{\partial V_O}{\partial V_I}\right|=1\]
\textbf{$V_{IL}$: druga oblast}\[\frac{V_{DD}-V_O}{R_L}=\frac{k_n}{2}\frac1{1+\frac{(V_I-V_{Tn})}{L_nE_{Cn}}}(V_I-V_{Tn})^2\]
Pretpostavka: $(V_{IL}-V_{Tn})\ll L_nE_{Cn}$.\[\frac{V_{DD}-V_O}{R_L}\approx\frac{k_n}{2}(V_I-V_{Tn})^2\]

\end{frame}

\slajdid{09_1-s019}
\begin{frame}[label=09_1-s019]{Prag VIL: implicitno diferenciranje}\setlength{\parskip}{2pt}\setlength{\abovedisplayskip}{4pt}\setlength{\belowdisplayskip}{4pt}\setlength{\abovedisplayshortskip}{3pt}\setlength{\belowdisplayshortskip}{3pt}
% element: 09_1-s019:e01
% element: 09_1-s019:e02
% element: 09_1-s019:e03
% element: 09_1-s019:e04
Diferenciramo obe strane po $V_I$:\[\frac{V_{DD}-V_O}{R_L}\approx\frac{k_n}{2}(V_I-V_{Tn})^2\quad\Bigg|\frac\partial{\partial V_I}\]

U tački praga: $\partial V_O/\partial V_I=-1$.
\[\frac1{R_L}\approx k_n(V_I-V_{Tn})\]
\begin{columns}[T,onlytextwidth]\column{.48\textwidth}\textbf{Ulazni napon}\[V_{IL}\approx V_{Tn}+\frac1{k_nR_L}\]
\column{.48\textwidth}\textbf{Pripadajući izlazni napon}\[V_{O(IL)}\approx V_{DD}-\frac1{2k_nR_L}\]
\end{columns}
\end{frame}

\slajdid{09_1-s020}
\begin{frame}[label=09_1-s020]{Prag VIH: triodna oblast}\setlength{\parskip}{2pt}\setlength{\abovedisplayskip}{4pt}\setlength{\belowdisplayskip}{4pt}\setlength{\abovedisplayshortskip}{3pt}\setlength{\belowdisplayshortskip}{3pt}
% element: 09_1-s020:e01
% element: 09_1-s020:e02
% element: 09_1-s020:e03
% element: 09_1-s020:e04
\textbf{$V_{IH}$: treća oblast}\[\frac{V_{DD}-V_O}{R_L}=\frac{k_n}{2}\frac1{1+\frac{V_O}{L_nE_{Cn}}}\bigl(2V_O(V_I-V_{Tn})-V_O^2\bigr)\]
Očekuje se nizak izlazni napon: $V_O<L_nE_{Cn}$.\[\frac{V_{DD}-V_O}{R_L}\approx\frac{k_n}{2}\bigl(2V_O(V_I-V_{Tn})-V_O^2\bigr)\]
Diferenciranje po $V_I$:\[-\frac1{R_L}\frac{\partial V_O}{\partial V_I}\approx\frac{k_n}{2}\left(2\frac{\partial V_O}{\partial V_I}(V_I-V_{Tn})+2V_O-2V_O\frac{\partial V_O}{\partial V_I}\right)\]

\end{frame}

\slajdid{09_1-s021}
\begin{frame}[label=09_1-s021]{Prag VIH: dve jednačine, dve nepoznate}\setlength{\parskip}{2pt}\setlength{\abovedisplayskip}{4pt}\setlength{\belowdisplayskip}{4pt}\setlength{\abovedisplayshortskip}{3pt}\setlength{\belowdisplayshortskip}{3pt}
% element: 09_1-s021:e01
% element: 09_1-s021:e02
% element: 09_1-s021:e03
% element: 09_1-s021:e04
\[\frac{\partial V_O}{\partial V_I}=-1\]

\textbf{Sistem za $V_{IH}$ i $V_{O(IH)}$}\[\color{DEcrvena}\left\{\color{DEtekst}\begin{aligned}
\alert{V_{IH}}&\approx2\alert{V_{O(IH)}}+V_{Tn}-\frac1{k_nR_L}\\[8pt]
\frac{V_{DD}-\alert{V_{O(IH)}}}{R_L}&\approx\frac{k_n}{2}\bigl(2\alert{V_{O(IH)}}(\alert{V_{IH}}-V_{Tn})-\alert{V_{O(IH)}^2}\bigr)
\end{aligned}\right.\]
\begin{columns}[T,onlytextwidth]\column{.48\textwidth}\[V_{O(IH)}\approx\sqrt{\frac{2V_{DD}}{3k_nR_L}}\]
\column{.48\textwidth}\[V_{IH}\approx V_{Tn}+\sqrt{\frac{8V_{DD}}{3k_nR_L}}-\frac1{k_nR_L}\]
\end{columns}
\end{frame}

\slajdid{09_1-s022}
\begin{frame}[label=09_1-s022]{Ispit: postupak analize}\setlength{\parskip}{2pt}\setlength{\abovedisplayskip}{4pt}\setlength{\belowdisplayskip}{4pt}\setlength{\abovedisplayshortskip}{3pt}\setlength{\belowdisplayshortskip}{3pt}
% element: 09_1-s022:e01
% element: 09_1-s022:e02
% element: 09_1-s022:e03
Uz odgovarajuće brojne vrednosti mogu se proveriti zanemarivanja.
\begin{enumerate}
\item Postavka preko „tačnih“ jednačina.
\item Računanje uz odgovarajuća zanemarivanja.
\item Održivost naponskih nivoa: da li $V_{OH}$ daje $V_{OL}$ i obrnuto; da li su $V_{IL}$ i $V_{IH}$ na odgovarajućim mestima.
\item Računanje i provera margina šuma.
\end{enumerate}
\medskip\alert{Na ispitu nije potrebna provera korektnosti prethodnih i sličnih zanemarivanja.}
\end{frame}

\slajdid{09_1-s023}
\begin{frame}[label=09_1-s023]{Održivost naponskih nivoa}\setlength{\parskip}{2pt}\setlength{\abovedisplayskip}{4pt}\setlength{\belowdisplayskip}{4pt}\setlength{\abovedisplayshortskip}{3pt}\setlength{\belowdisplayshortskip}{3pt}
% element: 09_1-s023:e01
% element: 09_1-s023:e02
% element: 09_1-s023:e03
% element: 09_1-s023:e04
% element: 09_1-s023:e05
\begin{columns}[T,onlytextwidth]\column{.28\textwidth}\centering \begin{tikzpicture}[x=.67cm,y=.67cm,font=\sitneoznake,>=Latex,remember picture,line width=.7pt]
\draw[->](-.13,0)--(3.45,0)node[right]{$V_I$};\draw[->](0,-.13)--(0,2.9)node[above]{$V_O$};\node[left]at(0,2.4){$V_{OH}$};\node[left]at(0,.13){$V_{OL}$};
\draw(0,2.4)--(.28,2.4)..controls(.6,2.4)and(.65,2.05)..(.83,1.65)..controls(1.15,.7)and(1.3,.4)..(1.6,.3)--(2.85,.13);
\draw[dashed](.58,-.1)node[below]{$V_{IL}$}--(.58,2.26);\draw[dashed](1.54,-.1)node[below]{$V_{IH}$}--(1.54,.36);\node[below]at(2.85,0){$V_{DD}$};\draw[dashed](0,.13)--(2.85,.13);
\draw[DEcrvena,dashed](0,2.26)--(.58,2.26);\draw[DEcrvena,dashed](0,.36)--(1.54,.36);
\draw(.4,2.44)--(.78,2.06);\node[right]at(.73,2.43){$a{=}{-}1$};\draw(1.31,.59)--(1.77,.13);\node[right]at(2.5,.85){$a{=}{-}1$};\coordinate(s23from)at(0,.36);
\draw[->,DEcrvena](2.2,1.5)--(1.54,.42);\node[above]at(2.05,1.65){$V_{O(IH)}$};\end{tikzpicture}
\column{.68\textwidth}\textbf{Logička nula mora zakočiti sledeći tranzistor}
\[\frac{V_{DD}}{k_nR_L(V_{DD}-V_{Tn})+1}<V_{Tn}\]
\[\alert{V_{O(IH)}<V_{IL}},\qquad\text{poželjno}\quad\alert{V_{O(IH)}<V_{Tn}}\]
\[\sqrt{\frac{2V_{DD}}{3k_nR_L}}<V_{Tn}\]
\end{columns}\medskip Za konstantno napajanje podešavamo $k_nR_L$, uz $\beta=k_nR_L/V_{DD}$.\[\sqrt{\frac2{3\beta}}<V_{Tn}\ \Rightarrow\ \beta>\frac2{3V_{Tn}^2}\]
Potreban je i uslov \alert{$V_{O(IL)}>V_{IH}$}.
\end{frame}

\slajdid{09_1-s024}
\begin{frame}[label=09_1-s024]{Prag odlučivanja: kvadratna jednačina}\setlength{\parskip}{2pt}\setlength{\abovedisplayskip}{4pt}\setlength{\belowdisplayskip}{4pt}\setlength{\abovedisplayshortskip}{3pt}\setlength{\belowdisplayshortskip}{3pt}
% element: 09_1-s024:e01
% element: 09_1-s024:e02
% element: 09_1-s024:e03
\textbf{Druga oblast; $V_I=V_O=V_S$}\[\frac{V_{DD}-V_S}{R_L}=\frac{k_n}{2}\frac{L_nE_{Cn}}{L_nE_{Cn}+(V_S-V_{Tn})}(V_S-V_{Tn})^2\]
Bez dodatnih zanemarivanja, rešavamo po $(V_S-V_{Tn})$:
\[(V_S-V_{Tn})^2+\frac{L_nE_{Cn}-(V_{DD}-V_{Tn})}{\frac{k_nR_LL_nE_{Cn}}2+1}(V_S-V_{Tn})-\frac{L_nE_{Cn}(V_{DD}-V_{Tn})}{\frac{k_nR_LL_nE_{Cn}}2+1}=0\]
\medskip Zamenimo brojne vrednosti i rešimo.
\end{frame}

\slajdid{09_1-s025}
\begin{frame}[label=09_1-s025]{Prag odlučivanja: približan račun}\setlength{\parskip}{2pt}\setlength{\abovedisplayskip}{4pt}\setlength{\belowdisplayskip}{4pt}\setlength{\abovedisplayshortskip}{3pt}\setlength{\belowdisplayshortskip}{3pt}
% element: 09_1-s025:e01
% element: 09_1-s025:e02
% element: 09_1-s025:e03
% element: 09_1-s025:e04
Očekujemo tačku blizu $V_{DD}/2$; tu vrednost unosimo u članove koji otežavaju račun.\[\frac{V_{DD}-\frac{V_{DD}}2}{R_L}=\frac{k_n}{2}\frac{L_nE_{Cn}}{L_nE_{Cn}+\left(\frac{V_{DD}}2-V_{Tn}\right)}(V_S-V_{Tn})^2\]
\[V_S-V_{Tn}=\sqrt{\frac{V_{DD}\left(L_nE_{Cn}+\left(\frac{V_{DD}}2-V_{Tn}\right)\right)}{k_nR_LL_nE_{Cn}}}=\alert{0.26V\ \Rightarrow\ V_S=0.71V}\]
\begin{itemize}
\item Tačan rezultat za standardne vrednosti: \alert{$V_S=0.7\,\mathrm V$}.
\item Polovina napajanja: $V_{DD}/2=0.6\,\mathrm V$.
\item Greška je mala: dobijeno je $V_S=0.71\,\mathrm V$.
\end{itemize}
\alert{Na ispitu je približan postupak dozvoljen i poželjan.}
\end{frame}

\slajdid{09_1-s026}
\begin{frame}[label=09_1-s026]{Ulazne struje i faktor grananja}\setlength{\parskip}{2pt}\setlength{\abovedisplayskip}{4pt}\setlength{\belowdisplayskip}{4pt}\setlength{\abovedisplayshortskip}{3pt}\setlength{\belowdisplayshortskip}{3pt}
% element: 09_1-s026:e01
% element: 09_1-s026:e02
% element: 09_1-s026:e03
Zanemarujući struje curenja:\[I_{IL}=I_{IH}\approx0\]
\begin{itemize}
\item Na osnovu statičkih struja faktor grananja bio bi beskonačan.
\item Ulazne kapacitivnosti MOSFET tranzistora veće su nego kod bipolarnih tranzistora.
\item Faktor grananja određuje dinamički režim MOS logičkih kola.
\end{itemize}
\end{frame}

\slajdid{09_1-s027}
\begin{frame}[label=09_1-s027]{Izlazni strujni kapaciteti}\setlength{\parskip}{2pt}\setlength{\abovedisplayskip}{4pt}\setlength{\belowdisplayskip}{4pt}\setlength{\abovedisplayshortskip}{3pt}\setlength{\belowdisplayshortskip}{3pt}
% element: 09_1-s027:e01
% element: 09_1-s027:e02
% element: 09_1-s027:e03
% element: 09_1-s027:e04
% element: 09_1-s027:e05
\textbf{Logička jedinica}\[I_{OH}=-\frac{V_{DD}-V_{OHmin}}{R_L}\]
\textbf{Logička nula — omska oblast}\[I_{OL}=\frac{k_n}{2}\frac1{1+\frac{V_{OLmax}}{L_nE_{Cn}}}\bigl(2V_{OLmax}(V_I-V_{Tn})-V_{OLmax}^2\bigr)-\frac{V_{DD}-V_{OLmax}}{R_L}\]
$V_{OHmin}=V_{IH}+\Delta$, $V_{OLmax}=V_{IL}-\Delta$: prostor za šum.\par Za najgori ulazni slučaj uzimamo $V_I=V_{OHmin}$:\[I_{OL}=\frac{k_n}{2}\frac1{1+\frac{V_{OLmax}}{L_nE_{Cn}}}\bigl(2V_{OLmax}(V_{OHmin}-V_{Tn})-V_{OLmax}^2\bigr)-\frac{V_{DD}-V_{OLmax}}{R_L}\]
\medskip\alert{Na ispitu, ako nije specificirano: $V_{OLmax}=V_{IL}$ i $V_{OHmin}=V_{IH}$.}
\end{frame}

\slajdid{09_1-s028}
\begin{frame}[label=09_1-s028]{Dinamički režim: punjenje i pražnjenje}\setlength{\parskip}{2pt}\setlength{\abovedisplayskip}{4pt}\setlength{\belowdisplayskip}{4pt}\setlength{\abovedisplayshortskip}{3pt}\setlength{\belowdisplayshortskip}{3pt}
% element: 09_1-s028:e01
% element: 09_1-s028:e02
% element: 09_1-s028:e03
% element: 09_1-s028:e04
\begin{columns}[T,onlytextwidth]\column{.48\textwidth}\textbf{Prelaz L u H}\par $\tau=R_LC_L$, $t_{pLH}=0.69\tau$, $t_r=2.2\tau$.\par\medskip
\centering {\RazaviSetup
\begin{circuitikz}[labeltext,american resistors,line width=.7pt]
\draw[wire](0,2)to[american resistor,l=$R_L$](0,.7)--(1.1,.7)to[C,l=$C_L$](1.1,-.65);
\coordinate(supply)at(0,2);\RazaviSupply{supply}{V_{DD}}
\coordinate(gnd)at(1.1,-.65);\RazaviGround{gnd}
\draw[wire](1.1,.7)node[junction]{}--(1.65,.7)node[rz terminal]{};
\draw[flow](2,-.6)--(2,.65);\node[right]at(2,.1){$v_o(t)$};\node[right]at(2,.65){$+$};
\coordinate(vg)at(2,-.7);\RazaviGround{vg}
\draw[flow](-.5,1.8)--(-.5,.9);\node[left]at(-.5,1.35){$i_{RC}(t)$};
\node[above]at(1.5,1.4){L u H};\end{circuitikz}}
\column{.48\textwidth}\textbf{Prelaz H u L}\par $T_D$ počinje u zasićenju; zanemarujemo $\lambda$.\[I_{Dn,D}=\frac{k_n}{2}\frac1{1+\frac{(V_I-V_{Tn})}{L_nE_{Cn}}}(V_I-V_{Tn})^2\]
Najgori slučaj: $V_I=V_{OHmin}$.\par\centering {\RazaviSetup
\begin{circuitikz}[labeltext,american resistors,line width=.7pt]
\draw[wire](0,2)to[american resistor,l=$R_L$](0,.7)--(1.1,.7)to[C,l=$C_L$](1.1,-.65);
\coordinate(supply)at(0,2);\RazaviSupply{supply}{V_{DD}}
\coordinate(gnd)at(1.1,-.65);\RazaviGround{gnd}
\draw[wire](1.1,.7)node[junction]{}--(1.65,.7)node[rz terminal]{};
\draw[flow](2,-.6)--(2,.65);\node[right]at(2,.1){$v_o(t)$};\node[right]at(2,.65){$+$};
\coordinate(vg)at(2,-.7);\RazaviGround{vg}
\draw[flow](-.5,1.8)--(-.5,.9);\node[left]at(-.5,1.35){$i_{RC}(t)$};
\draw[wire](0,.7)node[junction]{}--(-.75,.7)to[american current source,l_=$I_{Dn}$](-.75,-.65);\coordinate(sg)at(-.75,-.65);\RazaviGround{sg}
\draw[flow](.75,-.5)--(.75,.5);\node[left]at(.75,-.25){$i_C(t)$};\node[above]at(1.5,1.4){H u L};
\end{circuitikz}}
\end{columns}\par $v(t_0^+)=V_{OH}$.


\end{frame}

\slajdid{09_1-s029}
\begin{frame}[label=09_1-s029]{Uslov zasićenja pri pražnjenju}\setlength{\parskip}{2pt}\setlength{\abovedisplayskip}{4pt}\setlength{\belowdisplayskip}{4pt}\setlength{\abovedisplayshortskip}{3pt}\setlength{\belowdisplayshortskip}{3pt}
% element: 09_1-s029:e01
% element: 09_1-s029:e02
% element: 09_1-s029:e03
% element: 09_1-s029:e04
% element: 09_1-s029:e05
\textbf{Uslov zasićenja}\[V_{DSn,D}=V_O\geq\frac{(V_{GSn}-V_{Tn})}{1+\frac{(V_{GSn}-V_{Tn})}{L_nE_{Cn}}}=\frac{(V_I-V_{Tn})}{1+\frac{(V_I-V_{Tn})}{L_nE_{Cn}}}\]
\textbf{U trenutku kašnjenja $t_{pHL}$}\[v_O(t_{pHL})=\frac{V_{OH}+V_{OL}}2\geq\frac{V_{OHmin}-V_{Tn}}{1+\frac{(V_{OHmin}-V_{Tn})}{L_nE_{Cn}}}\]
\[\text{Tačno}\qquad\frac{V_{OH}+V_{OL}}2\geq\frac{L_nE_{Cn}(V_{OHmin}-V_{Tn})}{L_nE_{Cn}+(V_{OHmin}-V_{Tn})}\]
Zanemarivanje:\[\frac{V_{OH}+V_{OL}}2\geq L_nE_{Cn}\quad\text{ili}\quad\frac{V_{OH}+V_{OL}}2\geq V_{OHmin}-V_{Tn}\]
Za standardne vrednosti tranzistor radi u zasićenju za sve vreme procesa pražnjenja.
\end{frame}

\slajdid{09_1-s030}
\begin{frame}[label=09_1-s030]{Dinamička otpornost MOS tranzistora}\setlength{\parskip}{2pt}\setlength{\abovedisplayskip}{4pt}\setlength{\belowdisplayskip}{4pt}\setlength{\abovedisplayshortskip}{3pt}\setlength{\belowdisplayshortskip}{3pt}
% element: 09_1-s030:e01
% element: 09_1-s030:e02
% element: 09_1-s030:e03
% element: 09_1-s030:e04
% element: 09_1-s030:e05
MOS u aktivnoj oblasti zamenjujemo ekvivalentnom otpornošću.\par\alert{Uključuje se i faktor $\lambda$.}\begin{center}{\RazaviSetup
\begin{circuitikz}[labeltext,american resistors,line width=.7pt]
\node[rz nmos](d)at(0,0){};\draw[wire](d.D)--(0,.65)--(1,.65)to[C,l=$C_L$](1,-.70);\coordinate(cg)at(1,-.70);\RazaviGround{cg}
\draw[wire](d.S)--(0,-.70);\coordinate(gnd)at(0,-.70);\RazaviGround{gnd}
\draw[wire](d.G)--++(-.35,0)node[rz terminal]{}node[left=2pt]{$V_{GS}$};
\node[above]at(.5,.65){$V_O\;(V_{DD}\to V_{DD}/2)$};\node[junction]at(0,.65){};
\end{circuitikz}}
\hspace{12mm}\begin{tikzpicture}[x=.8cm,y=.6cm,font=\sitneoznake,>=Latex,line width=.7pt]
\draw[->](-.12,0)--(3.7,0)node[right]{$V_{DS}$};\draw[->](0,-.1)--(0,2.65)node[left]{$I_D$};
\draw(0,0)..controls(.4,1.85)and(.45,1.85)..(.8,1.93)--(3.25,2.38);\node[right]at(3.25,2.38){$V_{GS}$};
\draw[dashed](0,2.36)--(3.15,2.36)--(3.15,0)node[below]{$V_{DD}$};\draw[dashed](0,2.08)--(1.6,2.08)--(1.6,0)node[below]{$V_{DD}/2$};
\draw(3.09,2.3)--(3.21,2.42);\draw(3.09,2.42)--(3.21,2.3);\node[below right]at(3.15,2.36){1};\draw(1.54,2.02)--(1.66,2.14);\draw(1.54,2.14)--(1.66,2.02);\node[below left]at(1.6,2.08){2};\draw[->](2.7,2.28)--(2.3,2.20);
\end{tikzpicture}
\end{center}\[t_p=0.69R_nC_L\]
\begin{itemize}\setlength{\itemsep}{1pt}
\item Aproksimacija je namenjena kratkom kanalu sa malim $V_{DSsat}$.
\item Uslov: $V_{DSsat}<V_{DD}/2$, za $V_{Tn}\leq V_{GS}\leq V_{DD}$.

\end{itemize}
\end{frame}

\slajdid{09_1-s031}
\begin{frame}[label=09_1-s031]{Srednja otpornost iz dve radne tačke}\setlength{\parskip}{2pt}\setlength{\abovedisplayskip}{4pt}\setlength{\belowdisplayskip}{4pt}\setlength{\abovedisplayshortskip}{3pt}\setlength{\belowdisplayshortskip}{3pt}
% element: 09_1-s031:e01
% element: 09_1-s031:e02
% element: 09_1-s031:e03
% element: 09_1-s031:e04
\begin{center}{\RazaviSetup
\begin{circuitikz}[labeltext,american resistors,line width=.7pt]
\node[rz nmos](d)at(0,0){};\draw[wire](d.D)--(0,.65)--(1,.65)to[C,l=$C_L$](1,-.70);\coordinate(cg)at(1,-.70);\RazaviGround{cg}
\draw[wire](d.S)--(0,-.70);\coordinate(gnd)at(0,-.70);\RazaviGround{gnd}
\draw[wire](d.G)--++(-.35,0)node[rz terminal]{}node[left=2pt]{$V_{GS}$};
\node[above]at(.5,.65){$V_O\;(V_{DD}\to V_{DD}/2)$};\node[junction]at(0,.65){};
\end{circuitikz}}
\hspace{12mm}\begin{tikzpicture}[x=.8cm,y=.5cm,font=\sitneoznake,>=Latex,line width=.7pt]
\draw[->](-.12,0)--(3.7,0)node[right]{$V_{DS}$};\draw[->](0,-.1)--(0,2.65)node[left]{$I_D$};
\draw(0,0)..controls(.4,1.85)and(.45,1.85)..(.8,1.93)--(3.25,2.38);\node[right]at(3.25,2.38){$V_{GS}$};
\draw[dashed](0,2.36)--(3.15,2.36)--(3.15,0)node[below]{$V_{DD}$};\draw[dashed](0,2.08)--(1.6,2.08)--(1.6,0)node[below]{$V_{DD}/2$};
\draw(3.09,2.3)--(3.21,2.42);\draw(3.09,2.42)--(3.21,2.3);\node[below right]at(3.15,2.36){1};\draw(1.54,2.02)--(1.66,2.14);\draw(1.54,2.14)--(1.66,2.02);\node[below left]at(1.6,2.08){2};\draw[->](2.7,2.28)--(2.3,2.20);
\end{tikzpicture}
\end{center}\textbf{Prvi način: srednja vrednost u tačkama 1 i 2}\[R_n=\dfrac{R_{ON1}+R_{ON2}}2\]
\begin{columns}[T,onlytextwidth]\column{.48\textwidth}\[R_{ON1}=\frac{V_{DD}}{I_{Dn1}}\]
\column{.48\textwidth}\[R_{ON2}=\frac{V_{DD}/2}{I_{Dn2}}\]
\end{columns}
\end{frame}

\slajdid{09_1-s032}
\begin{frame}[label=09_1-s032]{Aproksimacija srednje otpornosti}\setlength{\parskip}{2pt}\setlength{\abovedisplayskip}{4pt}\setlength{\belowdisplayskip}{4pt}\setlength{\abovedisplayshortskip}{3pt}\setlength{\belowdisplayshortskip}{3pt}
% element: 09_1-s032:e01
% element: 09_1-s032:e02
% element: 09_1-s032:e03
% element: 09_1-s032:e04
Za skraćivanje uvodimo:\[I_{Dnsat}=\frac{k_n}{2}\frac{(V_{GS}-V_{Tn})^2}{1+\frac{V_{GS}-V_{Tn}}{L_nE_{Cn}}}\]
\[R_{ON1}=\frac{V_{DD}}{I_{Dn1}}=\frac{V_{DD}}{I_{Dnsat}(1+\lambda_nV_{DD})}\qquad R_{ON2}=\frac{\frac{V_{DD}}2}{I_{Dn2}}=\frac{V_{DD}}{2I_{Dnsat}(1+\lambda_n\frac{V_{DD}}2)}\]
\[R_n=\frac{V_{DD}}{4I_{Dnsat}}\left(\frac2{(1+\lambda_nV_{DD})}+\frac1{(1+\lambda_n\frac{V_{DD}}2)}\right)\]
Za $\lambda_nV_{DD}\ll1$: $1/(1+\lambda_nV_{DD})\approx1-\lambda_nV_{DD}$.\[R_n\approx\frac34\frac{V_{DD}}{I_{Dnsat}}\left(1-\frac56\lambda_nV_{DD}\right)\]
\medskip\alert{Koji $V_{GS}$ uzimamo za račun struje zasićenja?}
\end{frame}

\slajdid{09_1-s033}
\begin{frame}[label=09_1-s033]{Integralna srednja otpornost}\setlength{\parskip}{2pt}\setlength{\abovedisplayskip}{4pt}\setlength{\belowdisplayskip}{4pt}\setlength{\abovedisplayshortskip}{3pt}\setlength{\belowdisplayshortskip}{3pt}
% element: 09_1-s033:e01
% element: 09_1-s033:e02
% element: 09_1-s033:e03
% element: 09_1-s033:e04
% element: 09_1-s033:e05
\begin{columns}[T,onlytextwidth]\column{.48\textwidth}\textbf{Srednja vrednost po definiciji}\[R_n=\frac1{V_2-V_1}\int_{V_1}^{V_2}R_{ON}(V)dV\]
$R_{ON}(V)=V/I_{Dn}(V)$.
Oba načina su aproksimacije; integralni je tačniji.\par\[\frac56=0.833,\qquad\frac79=0.777\]
\column{.48\textwidth}Uz $I_{Dnsat}$:

\[R_n\approx\frac34\frac{V_{DD}}{I_{Dnsat}}\left(1-\frac79\lambda_nV_{DD}\right)\]
\end{columns}\[R_n=\frac1{\frac{V_{DD}}2-V_{DD}}\int_{V_{DD}}^{\frac{V_{DD}}2}\frac V{I_{Dnsat}(1+\lambda_nV)}dV\]

\end{frame}

\slajdid{09_1-s034}
\begin{frame}[label=09_1-s034]{Ekvivalentna otpornost PMOS tranzistora}\setlength{\parskip}{2pt}\setlength{\abovedisplayskip}{4pt}\setlength{\belowdisplayskip}{4pt}\setlength{\abovedisplayshortskip}{3pt}\setlength{\belowdisplayshortskip}{3pt}
% element: 09_1-s034:e01
% element: 09_1-s034:e02
% element: 09_1-s034:e03
% element: 09_1-s034:e04
PMOS u PUN mreži puni izlaznu kapacitivnost.\begin{center}{\RazaviSetup
\begin{circuitikz}[labeltext,american resistors,line width=.7pt]
\node[rz pmos](d)at(0,0){};\draw[wire](d.S)--(0,.7);\coordinate(supply)at(0,.7);\RazaviSupply{supply}{V_{DD}}
\draw[wire](d.D)--(0,-.65)--(1,-.65)to[C,l=$C_L$](1,-1.7);\coordinate(gnd)at(1,-1.7);\RazaviGround{gnd}
\draw[wire](d.G)--++(-.35,0)node[rz terminal]{}node[left=2pt]{$V_{SG}$};
\node[left]at(-.2,-.65){$V_O\;(0\to V_{DD}/2)$};\node[junction]at(0,-.65){};
\end{circuitikz}}
\hspace{12mm}\begin{tikzpicture}[x=.8cm,y=.6cm,font=\sitneoznake,>=Latex,line width=.7pt]
\draw[->](-.12,0)--(3.7,0)node[right]{$V_{SD}$};\draw[->](0,-.1)--(0,2.65)node[left]{$I_D$};
\draw(0,0)..controls(.4,1.85)and(.45,1.85)..(.8,1.93)--(3.25,2.38);\node[right]at(3.25,2.38){$V_{GS}$};
\draw[dashed](0,2.36)--(3.15,2.36)--(3.15,0)node[below]{$V_{DD}$};\draw[dashed](0,2.08)--(1.6,2.08)--(1.6,0)node[below]{$V_{DD}/2$};
\draw(3.09,2.3)--(3.21,2.42);\draw(3.09,2.42)--(3.21,2.3);\node[below right]at(3.15,2.36){1};\draw(1.54,2.02)--(1.66,2.14);\draw(1.54,2.14)--(1.66,2.02);\node[below left]at(1.6,2.08){2};\draw[->](2.7,2.28)--(2.3,2.20);
\end{tikzpicture}
\end{center}Sors je na $V_{DD}$, drejn na izlazu i kapacitivnosti.\[R_p\approx\frac34\frac{V_{DD}}{I_{Dpsat}}\left(1-\frac56|\lambda_p|V_{DD}\right)\]
\[R_p\approx\frac34\frac{V_{DD}}{I_{Dpsat}}\left(1-\frac79|\lambda_p|V_{DD}\right)\]

\end{frame}

\slajdid{09_1-s035}
\begin{frame}[label=09_1-s035]{Aktivno opterećenje: izbor parametara}\setlength{\parskip}{2pt}\setlength{\abovedisplayskip}{4pt}\setlength{\belowdisplayskip}{4pt}\setlength{\abovedisplayshortskip}{3pt}\setlength{\belowdisplayshortskip}{3pt}
% element: 09_1-s035:e01
% element: 09_1-s035:e02
% element: 09_1-s035:e03
% element: 09_1-s035:e04
\alert{NMOS invertor sa aktivnim opterećenjem: NMOS sa indukovanim kanalom.}\begin{enumerate}
\item Veći $R_L$: manji napon logičke nule i manji uticaj na njen strujni kapacitet.
\item Manji $R_L$: brže punjenje kapacitivnosti na izlazu.
\item $R_L$ utiče na $V_{IL}$ i $V_{IH}$: manji $R_L$ daje veće pragove i obrnuto.
\end{enumerate}
Potrebni su kompromis i izbor faktora $k_nR_L$.
\[k_n=\mu_nC_{oxn}\frac{\alert{W_n}}{L_n}\]
Za fiksne $\mu_n$, $C_{oxn}$ i $L_n$ kontrolišemo \alert{širinu kanala $W_n$}.
\end{frame}

\slajdid{09_1-s036}
\begin{frame}[label=09_1-s036]{NMOS opterećenje sa kontrolnim naponom}\setlength{\parskip}{2pt}\setlength{\abovedisplayskip}{4pt}\setlength{\belowdisplayskip}{4pt}\setlength{\abovedisplayshortskip}{3pt}\setlength{\belowdisplayshortskip}{3pt}
% element: 09_1-s036:e01
% element: 09_1-s036:e02
% element: 09_1-s036:e03
\begin{columns}[T,onlytextwidth]\column{.30\textwidth}\centering {\RazaviSetup
\begin{circuitikz}[labeltext,american resistors,line width=.7pt]
\node[rz nmos] (d) at (0,0) {};\node[rz nmos] (l) at (0,1.4) {};
\draw[wire](d.D)--(l.S);\coordinate(out)at(0,.70);
\draw[wire](out)node[junction]{}--(.55,.70)node[rz terminal]{}node[right=2pt]{$V_O$};
\draw[wire](l.D)--(0,2.10);\coordinate(supply)at(0,2.10);\RazaviSupply{supply}{V_{DD}}
\draw[wire](d.S)--(0,-.70);\coordinate(gnd)at(0,-.70);\RazaviGround{gnd}
\draw[wire](d.G)--++(-.35,0)node[rz terminal]{}node[left=2pt]{$V_I$};
\node[right=2pt]at(d.center){$T_D$};\node[right=2pt]at(l.center){$T_L$};
\draw[wire](l.G)--++(-.35,0)node[rz terminal]{}node[left=2pt]{$V_C$};\draw[flow](1.25,2.05)--(.48,1.55);
\end{circuitikz}}
\column{.66\textwidth}Otpornik $R_L$ zamenjujemo NMOS tranzistorom sa indukovanim kanalom.
\begin{itemize}
\item $T_L$: load; gejt je na kontrolnom naponu $V_C$.
\item Drejn $T_L$: $V_{DD}$; sors: izlazni napon.
\item Režim opterećenja određuje $V_C$.
\item Za kvalitativnu analizu koristimo \alert{dugačak kanal}.
\end{itemize}\end{columns}
\end{frame}

\slajdid{09_1-s037}
\begin{frame}[label=09_1-s037]{Režim rada tranzistora opterećenja}\setlength{\parskip}{2pt}\setlength{\abovedisplayskip}{4pt}\setlength{\belowdisplayskip}{4pt}\setlength{\abovedisplayshortskip}{3pt}\setlength{\belowdisplayshortskip}{3pt}
% element: 09_1-s037:e01
% element: 09_1-s037:e02
% element: 09_1-s037:e03
% element: 09_1-s037:e04
\begin{columns}[T,onlytextwidth]\column{.24\textwidth}\centering {\RazaviSetup
\begin{circuitikz}[labeltext,american resistors,line width=.7pt]
\node[rz nmos] (d) at (0,0) {};\node[rz nmos] (l) at (0,1.4) {};
\draw[wire](d.D)--(l.S);\coordinate(out)at(0,.70);
\draw[wire](out)node[junction]{}--(.55,.70)node[rz terminal]{}node[right=2pt]{$V_O$};
\draw[wire](l.D)--(0,2.10);\coordinate(supply)at(0,2.10);\RazaviSupply{supply}{V_{DD}}
\draw[wire](d.S)--(0,-.70);\coordinate(gnd)at(0,-.70);\RazaviGround{gnd}
\draw[wire](d.G)--++(-.35,0)node[rz terminal]{}node[left=2pt]{$V_I$};
\node[right=2pt]at(d.center){$T_D$};\node[right=2pt]at(l.center){$T_L$};
\draw[wire](l.G)--++(-.35,0)node[rz terminal]{}node[left=2pt]{$V_C$};
\end{circuitikz}}
\column{.72\textwidth}\textbf{Naponi na $T_L$}\[V_{GS,L}=V_{G,L}-V_{S,L}=V_C-V_O\]

\textbf{Uslov zasićenja}\[V_{DS,L}\geq V_{GS,L}-V_{Tn,L}\]

\[V_{DD}+V_{Tn,L}\geq V_C\]
Podešavanjem kontrolnog napona:\begin{columns}[T,onlytextwidth]\column{.48\textwidth}Zasićenje\[V_C<V_{DD}+V_{Tn,L}\]
\column{.48\textwidth}Omska oblast\[V_C>V_{DD}+V_{Tn,L}\]
\end{columns}\end{columns}
\end{frame}

\slajdid{09_1-s038}
\begin{frame}[label=09_1-s038]{Aktivno omsko opterećenje: logička jedinica}\setlength{\parskip}{2pt}\setlength{\abovedisplayskip}{4pt}\setlength{\belowdisplayskip}{4pt}\setlength{\abovedisplayshortskip}{3pt}\setlength{\belowdisplayshortskip}{3pt}
% element: 09_1-s038:e01
% element: 09_1-s038:e02
% element: 09_1-s038:e03
% element: 09_1-s038:e04
\begin{columns}[T,onlytextwidth]\column{.26\textwidth}\centering {\RazaviSetup
\begin{circuitikz}[labeltext,american resistors,line width=.7pt]
\node[rz nmos] (d) at (0,0) {};\node[rz nmos] (l) at (0,1.4) {};
\draw[wire](d.D)--(l.S);\coordinate(out)at(0,.70);
\draw[wire](out)node[junction]{}--(.55,.70)node[rz terminal]{}node[right=2pt]{$V_O$};
\draw[wire](l.D)--(0,2.10);\coordinate(supply)at(0,2.10);\RazaviSupply{supply}{V_{DD}}
\draw[wire](d.S)--(0,-.70);\coordinate(gnd)at(0,-.70);\RazaviGround{gnd}
\draw[wire](d.G)--++(-.35,0)node[rz terminal]{}node[left=2pt]{$V_I$};
\node[right=2pt]at(d.center){$T_D$};\node[right=2pt]at(l.center){$T_L$};
\draw[wire](l.G)--++(-.35,0)node[rz terminal]{}node[left=2pt]{$V_C$};
\end{circuitikz}}
\column{.70\textwidth}\textbf{Prvi slučaj: $T_L$ u omskoj oblasti}
\par Očekujemo rezultate slične pasivnom opterećenju.
\medskip\par $V_I<V_{Tn,D}$: pogonski tranzistor $T_D$ je zakočen.\par Neopterećeno kolo: $I_{Dn,D}=I_{Dn,L}=0$.
\[I_{Dn,L}=\frac{k_{n,L}}2\bigl(2V_{DS,L}(V_{GS,L}-V_{Tn,L})-V_{DS,L}^2\bigr)=0\]
Radna tačka: $V_{DS,L}=0$.\[V_O=V_{OH}=V_{DD}-V_{DS,L}=V_{DD}\]
\end{columns}
\end{frame}

\slajdid{09_1-s039}
\begin{frame}[label=09_1-s039]{Početak provođenja pogonskog tranzistora}\setlength{\parskip}{2pt}\setlength{\abovedisplayskip}{4pt}\setlength{\belowdisplayskip}{4pt}\setlength{\abovedisplayshortskip}{3pt}\setlength{\belowdisplayshortskip}{3pt}
% element: 09_1-s039:e01
% element: 09_1-s039:e02
% element: 09_1-s039:e03
% element: 09_1-s039:e04
\begin{columns}[T,onlytextwidth]\column{.25\textwidth}\centering {\RazaviSetup
\begin{circuitikz}[labeltext,american resistors,line width=.7pt]
\node[rz nmos] (d) at (0,0) {};\node[rz nmos] (l) at (0,1.4) {};
\draw[wire](d.D)--(l.S);\coordinate(out)at(0,.70);
\draw[wire](out)node[junction]{}--(.55,.70)node[rz terminal]{}node[right=2pt]{$V_O$};
\draw[wire](l.D)--(0,2.10);\coordinate(supply)at(0,2.10);\RazaviSupply{supply}{V_{DD}}
\draw[wire](d.S)--(0,-.70);\coordinate(gnd)at(0,-.70);\RazaviGround{gnd}
\draw[wire](d.G)--++(-.35,0)node[rz terminal]{}node[left=2pt]{$V_I$};
\node[right=2pt]at(d.center){$T_D$};\node[right=2pt]at(l.center){$T_L$};
\draw[wire](l.G)--++(-.35,0)node[rz terminal]{}node[left=2pt]{$V_C$};
\end{circuitikz}}
\column{.71\textwidth}$V_I=V_{Tn,D}+\varepsilon$: $T_D$ provodi malu struju.\[V_{DS,D}\geq V_{GS,D}-V_{Tn,D},\qquad V_O\geq V_I-V_{Tn,D}\]
\textbf{Ravnoteža $I_{Dn,L}=I_{Dn,D}$}\end{columns}\[\frac{k_{n,L}}2\bigl(2V_{DS,L}(V_{GS,L}-V_{Tn,L})-V_{DS,L}^2\bigr)=\frac{k_{n,D}}2(V_{GS,D}-V_{Tn,D})^2\]
\[\frac{k_{n,L}}2\bigl(2(V_{DD}-V_O)(V_C-V_O-V_{Tn,L})-(V_{DD}-V_O)^2\bigr)=\frac{k_{n,D}}2(V_I-V_{Tn,D})^2\]
\[k_{n,L}(V_{DD}-V_O)(2V_C-2V_{Tn,L}-V_{DD}-V_O)=k_{n,D}(V_I-V_{Tn,D})^2\]

\end{frame}

\slajdid{09_1-s040}
\begin{frame}[label=09_1-s040]{Prag VIL sa aktivnim omskim opterećenjem}\setlength{\parskip}{2pt}\setlength{\abovedisplayskip}{4pt}\setlength{\belowdisplayskip}{4pt}\setlength{\abovedisplayshortskip}{3pt}\setlength{\belowdisplayshortskip}{3pt}
% element: 09_1-s040:e01
% element: 09_1-s040:e02
% element: 09_1-s040:e03
% element: 09_1-s040:e04
% element: 09_1-s040:e05
$V_C=V_{DD}+V_{Tn,L}+\Delta$, za $V_C>V_{DD}+V_{Tn,L}$.\[\alert{1.}\quad k_{n,L}(V_{DD}-V_O)(V_{DD}+2\Delta-V_O)=k_{n,D}(V_I-V_{Tn,D})^2\]
Diferenciranje uz $\partial V_O/\partial V_I=-1$:\[k_{n,L}(V_{DD}+2\Delta-V_O)+k_{n,L}(V_{DD}-V_O)=2k_{n,D}(V_I-V_{Tn,D})\]
\[k_{n,L}(V_{DD}+\Delta-V_O)=k_{n,D}(V_I-V_{Tn,D})\]
\[\alert{2.}\quad V_I=V_{Tn,D}+\frac{k_{n,L}}{k_{n,D}}(V_{DD}+\Delta-V_O)\]
Zamenom u ravnotežu struja:\[k_{n,D}(V_{DD}-V_O)(V_{DD}+2\Delta-V_O)=k_{n,L}(V_{DD}+\Delta-V_O)^2\]

\end{frame}

\slajdid{09_1-s041}
\begin{frame}[label=09_1-s041]{Poseban slučaj i kvadratna jednačina}\setlength{\parskip}{2pt}\setlength{\abovedisplayskip}{4pt}\setlength{\belowdisplayskip}{4pt}\setlength{\abovedisplayshortskip}{3pt}\setlength{\belowdisplayshortskip}{3pt}
% element: 09_1-s041:e01
% element: 09_1-s041:e02
% element: 09_1-s041:e03
% element: 09_1-s041:e04
\[k_{n,D}(V_{DD}-V_O)(V_{DD}+2\Delta-V_O)=k_{n,L}(V_{DD}+\Delta-V_O)^2\]
Za $\Delta=0$ i $k_{n,L}\ne k_{n,D}$:\[V_O=V_{DD},\qquad V_{IL}=V_{Tn,D}\]
Sređivanje po $V_{DD}-V_O$:\[k_{n,D}(V_{DD}-V_O)^2+2\Delta k_{n,D}(V_{DD}-V_O)=k_{n,L}\bigl((V_{DD}-V_O)^2+2\Delta(V_{DD}-V_O)+\Delta^2\bigr)\]
\[(V_{DD}-V_O)^2(k_{n,D}-k_{n,L})+2\Delta(V_{DD}-V_O)(k_{n,D}-k_{n,L})-\Delta^2 k_{n,L}=0\]
Za $k_{n,L}=k_{n,D}$ jednačina nema rešenje: nema tačke jediničnog pojačanja u ovoj oblasti.
\end{frame}

\slajdid{09_1-s042}
\begin{frame}[label=09_1-s042]{Rešenje za prag VIL}\setlength{\parskip}{2pt}\setlength{\abovedisplayskip}{4pt}\setlength{\belowdisplayskip}{4pt}\setlength{\abovedisplayshortskip}{3pt}\setlength{\belowdisplayshortskip}{3pt}
% element: 09_1-s042:e01
% element: 09_1-s042:e02
% element: 09_1-s042:e03
\textbf{Izlazni napon u tački praga}\[\begin{aligned}
V_{O(IL)}&=V_{DD}+(V_C-V_{DD}-V_{Tn,L})\left(1-\sqrt{\frac{k_{n,D}}{(k_{n,D}-k_{n,L})}}\right)\\[1pt]
&=V_{DD}+(V_C-V_{DD}-V_{Tn,L})\left(1-\sqrt{\frac1{1-k_{n,L}/k_{n,D}}}\right)
\end{aligned}\]
\smallskip\textbf{Ulazni prag}\[\begin{aligned}
V_{IL}&=V_{Tn,D}+(V_C-V_{DD}-V_{Tn.L})\frac{k_{n,L}}{k_{n,D}}\sqrt{\frac{k_{n,D}}{(k_{n,D}-k_{n,L})}}\\[1pt]
&=V_{Tn,D}+(V_C-V_{DD}-V_{Tn.L})\frac{k_{n,L}}{k_{n,D}}\sqrt{\frac1{1-k_{n,L}/k_{n,D}}}
\end{aligned}\]

\end{frame}

\slajdid{09_1-s043}
\begin{frame}[label=09_1-s043]{Oba tranzistora u omskoj oblasti}\setlength{\parskip}{2pt}\setlength{\abovedisplayskip}{4pt}\setlength{\belowdisplayskip}{4pt}\setlength{\abovedisplayshortskip}{3pt}\setlength{\belowdisplayshortskip}{3pt}
% element: 09_1-s043:e01
% element: 09_1-s043:e02
% element: 09_1-s043:e03
% element: 09_1-s043:e04
Napon na drejnu tranzistora T1 opada. Granica:\[V_O=V_{DS,D}=V_{GS,D}-V_{Tn,D}=V_I-V_{Tn,D}\]
\textbf{Ravnoteža struja}\[\frac{k_{n,L}}2\bigl(2V_{DS,L}(V_{GS,L}-V_{Tn,L})-V_{DS,L}^2\bigr)=\frac{k_{n,D}}2\bigl(2V_{DS,D}(V_{GS,D}-V_{Tn,D})-V_{DS,D}^2\bigr)\]
Uz $V_C=V_{DD}+V_{Tn,L}+\Delta$:\[\alert{1.}\quad k_{n,L}(V_{DD}-V_O)(V_{DD}+2\Delta-V_O)=k_{n,D}V_O(2V_I-2V_{Tn,D}-V_O)\]

\end{frame}

\slajdid{09_1-s044}
\begin{frame}[label=09_1-s044]{Prag VIH sa omskim opterećenjem}\setlength{\parskip}{2pt}\setlength{\abovedisplayskip}{4pt}\setlength{\belowdisplayskip}{4pt}\setlength{\abovedisplayshortskip}{3pt}\setlength{\belowdisplayshortskip}{3pt}
% element: 09_1-s044:e01
% element: 09_1-s044:e02
% element: 09_1-s044:e03
% element: 09_1-s044:e04
% element: 09_1-s044:e05
\[\alert{1.}\quad k_{n,L}(V_{DD}-V_O)(V_{DD}+2\Delta-V_O)=k_{n,D}V_O(2V_I-2V_{Tn,D}-V_O)\]
U tački $V_{IH}$: $\partial V_O/\partial V_I=-1$.\[k_{n,L}(V_{DD}+2\Delta-V_O)+k_{n,L}(V_{DD}-V_O)=-k_{n,D}(2V_I-2V_{Tn,D}-V_O)+k_{n,D}V_O(2+1)\]
\[k_{n,L}(2V_{DD}+2\Delta-2V_O)=k_{n,D}(2V_O-2V_I+2V_{Tn,D})\]
\[\alert{2.}\quad V_I=V_{Tn,D}+V_O-\frac{k_{n,L}}{k_{n,D}}(V_{DD}+\Delta-V_O)\]
Zamena u ravnotežu struja:\[k_{n,L}(V_{DD}-V_O)(V_{DD}+2\Delta-V_O)=k_{n,D}V_O\left(V_O-2\frac{k_{n,L}}{k_{n.D}}(V_{DD}+\Delta-V_O)\right)\]

\end{frame}

\slajdid{09_1-s045}
\begin{frame}[label=09_1-s045]{Rezultati za prag VIH}\setlength{\parskip}{2pt}\setlength{\abovedisplayskip}{4pt}\setlength{\belowdisplayskip}{4pt}\setlength{\abovedisplayshortskip}{3pt}\setlength{\belowdisplayshortskip}{3pt}
% element: 09_1-s045:e01
% element: 09_1-s045:e02
% element: 09_1-s045:e03
\textbf{Izlazni napon u tački praga}\[V_{O(IH)}\approx\frac{k_{n,L}}{k_{n,D}}(V_C-V_{Tn,L})\left(1+\sqrt{1+\frac{k_{n,D}}{k_{n,L}}\left(1-\frac{(V_C-V_{DD}-V_{Tn,L})^2}{(V_C-V_{Tn,L})^2}\right)}\right)\]
\bigskip\textbf{Ulazni prag}\[V_{IH}\approx V_{Tn,D}+\frac{k_{n,L}}{k_{n,D}}(V_C-V_{Tn,L})\sqrt{1+\frac{k_{n,D}}{k_{n,L}}\left(1-\frac{(V_C-V_{DD}-V_{Tn,L})^2}{(V_C-V_{Tn,L})^2}\right)}\]

\end{frame}

\slajdid{09_1-s046}
\begin{frame}[label=09_1-s046]{Logička nula sa omskim opterećenjem}\setlength{\parskip}{2pt}\setlength{\abovedisplayskip}{4pt}\setlength{\belowdisplayskip}{4pt}\setlength{\abovedisplayshortskip}{3pt}\setlength{\belowdisplayshortskip}{3pt}
% element: 09_1-s046:e01
% element: 09_1-s046:e02
% element: 09_1-s046:e03
% element: 09_1-s046:e04
Ravnoteža struja:\[k_{n,L}(V_{DD}-V_O)(V_{DD}+2\Delta-V_O)=k_{n,D}V_O(2V_I-2V_{Tn,D}-V_O)\]
Za $V_I=V_{OH}$:\[k_{n,L}(V_{DD}-V_O)(V_{DD}+2\Delta-V_O)=k_{n,D}V_O(2V_{OH}-2V_{Tn,D}-V_O)\]
Uz mali $V_{OL}$:\[k_{n,L}(V_{DD}-V_O)(V_{DD}+2\Delta)=k_{n,D}V_O(2V_{OH}-2V_{Tn,D})\]
\[V_{OL}\approx\frac12\frac{k_{n,L}}{k_{n,D}}\frac{V_{DD}(V_{DD}+2\Delta)}{V_{DD}-V_{Tn,D}}=\frac12\frac{k_{n,L}}{k_{n,D}}\frac{V_{DD}(2V_C-V_{DD}-V_{Tn,L})}{V_{DD}-V_{Tn,D}}\]

\end{frame}

\slajdid{09_1-s047}
\begin{frame}[label=09_1-s047]{Odnos dimenzija tranzistora}\setlength{\parskip}{2pt}\setlength{\abovedisplayskip}{4pt}\setlength{\belowdisplayskip}{4pt}\setlength{\abovedisplayshortskip}{3pt}\setlength{\belowdisplayshortskip}{3pt}
% element: 09_1-s047:e01
% element: 09_1-s047:e02
% element: 09_1-s047:e03
\[V_{OL}\approx\frac12\alert{\frac{k_{n,L}}{k_{n,D}}}\frac{V_{DD}(V_{DD}+2\Delta)}{V_{DD}-V_{Tn,D}}=\frac12\alert{\frac{k_{n,L}}{k_{n,D}}}\frac{V_{DD}(2V_C-V_{DD}-V_{Tn,L})}{V_{DD}-V_{Tn,D}}\]
\medskip\par\alert{Povećanje obe širine istim faktorom ne menja odnos parametara.}\[\frac{k_{n,L}}{k_{n,D}}=\frac{\mu_{n,L}C_{oxn,L}\dfrac{W_{n,L}}{L_{n,L}}}{\mu_{n,D}C_{oxn,D}\dfrac{W_{n,D}}{L_{n,D}}}=\frac{\mu_{n,L}C_{oxn,L}\dfrac{F\times W_{n,L}}{L_{n,L}}}{\mu_{n,D}C_{oxn,D}\dfrac{F\times W_{n,D}}{L_{n,D}}}\]

\end{frame}

\slajdid{09_1-s048}
\begin{frame}[label=09_1-s048]{Karakteristika prenosa i kompromisi}\setlength{\parskip}{2pt}\setlength{\abovedisplayskip}{4pt}\setlength{\belowdisplayskip}{4pt}\setlength{\abovedisplayshortskip}{3pt}\setlength{\belowdisplayshortskip}{3pt}
% element: 09_1-s048:e01
% element: 09_1-s048:e02
% element: 09_1-s048:e03
% element: 09_1-s048:e04
% element: 09_1-s048:e05
\begin{columns}[T,onlytextwidth]\column{.32\textwidth}\centering \begin{tikzpicture}[x=.67cm,y=.67cm,font=\sitneoznake,>=Latex,baseline=(current bounding box.center)]
\draw[->](-.13,0)--(3.45,0)node[right]{$V_I$};\draw[->](0,-.13)--(0,2.9)node[above]{$V_O$};\node[left]at(0,2.4){$V_{OH}$};\node[left]at(0,.13){$V_{OL}$};
\draw(0,2.4)--(.28,2.4)..controls(.6,2.4)and(.65,2.05)..(.83,1.65)..controls(1.15,.7)and(1.3,.4)..(1.6,.3)--(2.85,.13);
\draw[dashed](.58,-.1)node[below]{$V_{IL}$}--(.58,2.26);\draw[dashed](1.54,-.1)node[below]{$V_{IH}$}--(1.54,.36);\node[below]at(2.85,0){$V_{DD}$};\draw[dashed](0,.13)--(2.85,.13);

\draw(.4,2.44)--(.78,2.06);\node[right]at(.73,2.43){$a{=}{-}1$};\draw(1.31,.59)--(1.77,.13);\node[right]at(1.66,.9){$a{=}{-}1$};
\end{tikzpicture}

\column{.64\textwidth}Parametre podešavamo odnosom $k_{n,D}/k_{n,L}$, odnosno odnosom $W_{n,D}/W_{n,L}$.
\[r_{DSn,L}\approx\frac{1}{k_{n,L}(V_{GS,L}-V_{Tn,L})}\]
\begin{itemize}
\item Mali $V_I$, visok $V_O$: velika dinamička otpornost i veliko pojačanje.
\item Veliki $V_I$, mali $V_O$: mala otpornost, brzo punjenje kapacitivnosti.
\item Kompromis: veći $V_{OL}$ i manji strujni kapacitet logičke nule.
\end{itemize}\end{columns}
\end{frame}

\slajdid{09_1-s049}
\begin{frame}[label=09_1-s049]{Gejt opterećenja na napajanju}\setlength{\parskip}{2pt}\setlength{\abovedisplayskip}{4pt}\setlength{\belowdisplayskip}{4pt}\setlength{\abovedisplayshortskip}{3pt}\setlength{\belowdisplayshortskip}{3pt}
% element: 09_1-s049:e01
% element: 09_1-s049:e02
% element: 09_1-s049:e03
% element: 09_1-s049:e04
% element: 09_1-s049:e05
\begin{columns}[T,onlytextwidth]\column{.26\textwidth}\centering {\RazaviSetup
\begin{circuitikz}[labeltext,american resistors,line width=.7pt]
\node[rz nmos] (d) at (0,0) {};\node[rz nmos] (l) at (0,1.4) {};
\draw[wire](d.D)--(l.S);\coordinate(out)at(0,.70);
\draw[wire](out)node[junction]{}--(.55,.70)node[rz terminal]{}node[right=2pt]{$V_O$};
\draw[wire](l.D)--(0,2.10);\coordinate(supply)at(0,2.10);\RazaviSupply{supply}{V_{DD}}
\draw[wire](d.S)--(0,-.70);\coordinate(gnd)at(0,-.70);\RazaviGround{gnd}
\draw[wire](d.G)--++(-.35,0)node[rz terminal]{}node[left=2pt]{$V_I$};
\node[right=2pt]at(d.center){$T_D$};\node[right=2pt]at(l.center){$T_L$};
\draw[wire](l.G)--++(-.25,0)|-(supply)node[junction]{};
\end{circuitikz}}
\column{.70\textwidth}Izbegavamo dodatni izvor: gejt $T_L$ povezujemo na $V_{DD}$.\[V_{GS,L}=V_{G,L}-V_{S,L}=V_{DD}-V_O\]
\[V_{DS,L}=V_{D,L}-V_{S,L}=V_{DD}-V_O\]
\textbf{Uslov zasićenja}\[V_{DS,L}\geq V_{GS,L}-V_{Tn,L}\]
\[V_{DD}-V_O\geq V_{DD}-V_O-V_{Tn,L}\]
Tranzistor T2 uvek radi u zasićenju. Važi i za kratki kanal.\end{columns}
\end{frame}

\slajdid{09_1-s050}
\begin{frame}[label=09_1-s050]{Logička jedinica sa zasićenim opterećenjem}\setlength{\parskip}{2pt}\setlength{\abovedisplayskip}{4pt}\setlength{\belowdisplayskip}{4pt}\setlength{\abovedisplayshortskip}{3pt}\setlength{\belowdisplayshortskip}{3pt}
% element: 09_1-s050:e01
% element: 09_1-s050:e02
% element: 09_1-s050:e03
% element: 09_1-s050:e04
\begin{columns}[T,onlytextwidth]\column{.27\textwidth}\centering {\RazaviSetup
\begin{circuitikz}[labeltext,american resistors,line width=.7pt]
\node[rz nmos] (d) at (0,0) {};\node[rz nmos] (l) at (0,1.4) {};
\draw[wire](d.D)--(l.S);\coordinate(out)at(0,.70);
\draw[wire](out)node[junction]{}--(.55,.70)node[rz terminal]{}node[right=2pt]{$V_O$};
\draw[wire](l.D)--(0,2.10);\coordinate(supply)at(0,2.10);\RazaviSupply{supply}{V_{DD}}
\draw[wire](d.S)--(0,-.70);\coordinate(gnd)at(0,-.70);\RazaviGround{gnd}
\draw[wire](d.G)--++(-.35,0)node[rz terminal]{}node[left=2pt]{$V_I$};
\node[right=2pt]at(d.center){$T_D$};\node[right=2pt]at(l.center){$T_L$};
\draw[wire](l.G)--++(-.25,0)|-(supply)node[junction]{};
\end{circuitikz}}
\column{.69\textwidth}$V_I<V_{Tn,D}$: $T_D$ je zakočen.\par Neopterećeno kolo: $I_{Dn,D}=I_{Dn,L}=0$.\[I_{Dn,L}=\frac{k_{n,L}}2(V_{GS,L}-V_{Tn,L})^2=0\]
Radna tačka: $V_{GS,L}=V_{Tn,L}$.\[V_O=V_{OH}=V_{DD}-V_{GS,L}=V_{DD}-V_{Tn,L}\]
\alert{Logička jedinica je niža nego kod omskog opterećenja.}\end{columns}
\end{frame}

\slajdid{09_1-s051}
\begin{frame}[label=09_1-s051]{Oba tranzistora u zasićenju}\setlength{\parskip}{2pt}\setlength{\abovedisplayskip}{4pt}\setlength{\belowdisplayskip}{4pt}\setlength{\abovedisplayshortskip}{3pt}\setlength{\belowdisplayshortskip}{3pt}
% element: 09_1-s051:e01
% element: 09_1-s051:e02
% element: 09_1-s051:e03
% element: 09_1-s051:e04
% element: 09_1-s051:e05
\begin{columns}[T,onlytextwidth]\column{.23\textwidth}\centering {\RazaviSetup
\begin{circuitikz}[labeltext,american resistors,line width=.7pt]
\node[rz nmos] (d) at (0,0) {};\node[rz nmos] (l) at (0,1.4) {};
\draw[wire](d.D)--(l.S);\coordinate(out)at(0,.70);
\draw[wire](out)node[junction]{}--(.55,.70)node[rz terminal]{}node[right=2pt]{$V_O$};
\draw[wire](l.D)--(0,2.10);\coordinate(supply)at(0,2.10);\RazaviSupply{supply}{V_{DD}}
\draw[wire](d.S)--(0,-.70);\coordinate(gnd)at(0,-.70);\RazaviGround{gnd}
\draw[wire](d.G)--++(-.35,0)node[rz terminal]{}node[left=2pt]{$V_I$};
\node[right=2pt]at(d.center){$T_D$};\node[right=2pt]at(l.center){$T_L$};
\draw[wire](l.G)--++(-.25,0)|-(supply)node[junction]{};
\end{circuitikz}}
\column{.73\textwidth}$V_I=V_{Tn,D}+\varepsilon$; $T_D$ provodi malu struju.\par $V_{DS,D}\geq V_{GS,D}-V_{Tn,D1}$, odnosno $V_O\geq V_I-V_{Tn,D}$.
\[\alert{1.}\quad k_{n,L}(V_{DD}-V_O-V_{Tn,L})^2=k_{n,D}(V_I-V_{Tn,D})^2\]\end{columns}
Za $\partial V_O/\partial V_I=-1$:
\[\alert{2.}\quad V_I=V_{Tn,D}+\frac{k_{n,L}}{k_{n,D}}(V_{DD}-V_O-V_{Tn,L})\]
Zamena u početnu jednačinu:\[k_{n,L}(V_{DD}-V_O-V_{Tn,L})^2=k_{n,L}(V_{DD}-V_O-V_{Tn,L})^2\quad\text{????}\]

\end{frame}

\slajdid{09_1-s052}
\begin{frame}[label=09_1-s052]{Konstantno pojačanje u oblasti zasićenja}\setlength{\parskip}{2pt}\setlength{\abovedisplayskip}{4pt}\setlength{\belowdisplayskip}{4pt}\setlength{\abovedisplayshortskip}{3pt}\setlength{\belowdisplayshortskip}{3pt}
% element: 09_1-s052:e01
% element: 09_1-s052:e02
% element: 09_1-s052:e03
% element: 09_1-s052:e04
% element: 09_1-s052:e05
U ovoj oblasti nema tačke sa pojačanjem $-1$; pojačanje je konstantno.\[k_{n,L}(V_{DD}-V_O-V_{Tn,L})^2=k_{n,D}(V_I-V_{Tn,D})^2\]
Eksplicitna karakteristika:\[V_O=V_{DD}-V_{Tn,L}-\sqrt{\frac{k_{n,D}}{k_{n,L}}}(V_I-V_{Tn,D})\]
\[a=\frac{dV_O}{dV_I}=-\sqrt{\frac{k_{n,D}}{k_{n,L}}}\]
Za logičko kolo:\begin{columns}[T,onlytextwidth]\column{.48\textwidth}\[\sqrt{\frac{k_{n,D}}{k_{n,L}}}>1\]
\column{.48\textwidth}\[k_{n,D}>k_{n,L}\]
\end{columns}
\end{frame}

\slajdid{09_1-s053}
\begin{frame}[label=09_1-s053]{Prelazak u omsku oblast i prag VIH}\setlength{\parskip}{2pt}\setlength{\abovedisplayskip}{4pt}\setlength{\belowdisplayskip}{4pt}\setlength{\abovedisplayshortskip}{3pt}\setlength{\belowdisplayshortskip}{3pt}
% element: 09_1-s053:e01
% element: 09_1-s053:e02
% element: 09_1-s053:e03
% element: 09_1-s053:e04
% element: 09_1-s053:e05
\[V_{IL}=V_{Tn,D}\]
Granica oblasti pogonskog tranzistora:\[V_O=V_{DS,D}=V_{GS,D}-V_{Tn,D}=V_I-V_{Tn,D}\]
Ravnoteža struja:\[\alert{1.}\quad k_{n,L}(V_{DD}-V_O-V_{Tn,L})^2=k_{n,D}V_O(2V_I-2V_{Tn,D}-V_O)\]
Za $\partial V_O/\partial V_I=-1$:
\[2k_{n,L}(V_{DD}-V_O-V_{Tn,L})=k_{n,D}(2V_O-2V_I+2V_{Tn,D})\]
\[\alert{2.}\quad V_I=V_{Tn,D}+V_O-\frac{k_{n,L}}{k_{n,D}}(V_{DD}-V_O-V_{Tn,L})\]

\end{frame}

\slajdid{09_1-s054}
\begin{frame}[label=09_1-s054]{Rezultati za zasićeno opterećenje}\setlength{\parskip}{2pt}\setlength{\abovedisplayskip}{4pt}\setlength{\belowdisplayskip}{4pt}\setlength{\abovedisplayshortskip}{3pt}\setlength{\belowdisplayshortskip}{3pt}
% element: 09_1-s054:e01
% element: 09_1-s054:e02
\textbf{Izlazni napon u tački praga}\[V_{O(IH)}\approx\frac{k_{n,L}}{k_{n,D}}(V_{DD}-V_{Tn,L})\left(1+\sqrt{1+\frac{k_{n,D}}{k_{n,L}}}\right)\]
\bigskip\textbf{Ulazni prag}\[V_{IH}\approx V_{Tn,D}+\frac{k_{n,L}}{k_{n,D}}(V_{DD}-V_{Tn,L})\sqrt{1+\frac{k_{n,D}}{k_{n,L}}}\]

\end{frame}

\slajdid{09_1-s055}
\begin{frame}[label=09_1-s055]{Logička nula sa zasićenim opterećenjem}\setlength{\parskip}{2pt}\setlength{\abovedisplayskip}{4pt}\setlength{\belowdisplayskip}{4pt}\setlength{\abovedisplayshortskip}{3pt}\setlength{\belowdisplayshortskip}{3pt}
% element: 09_1-s055:e01
% element: 09_1-s055:e02
% element: 09_1-s055:e03
% element: 09_1-s055:e04
% element: 09_1-s055:e05
\[k_{n,L}(V_{DD}-V_O-V_{Tn,L})^2=k_{n,D}V_O(2V_I-2V_{Tn,D}-V_O)\]
\begin{columns}[T,onlytextwidth]\column{.65\textwidth}Za $V_I=V_{OH}$:\[k_{n,L}(V_{DD}-V_O-V_{Tn,L})^2=k_{n,D}V_O(2V_{OH}-2V_{Tn,D}-V_O)\]
Uz mali $V_{OL}$:\[k_{n,L}(V_{DD}-V_{Tn,L})^2=k_{n,D}V_O(2V_{OH}-2V_{Tn,D})\]
\column{.31\textwidth}\centering \begin{tikzpicture}[x=.67cm,y=.67cm,font=\sitneoznake,>=Latex,baseline=(current bounding box.center)]
\draw[->](-.13,0)--(3.45,0)node[right]{$V_I$};\draw[->](0,-.13)--(0,2.9)node[above]{$V_O$};\node[left]at(0,2.4){$V_{OH}$};\node[left]at(0,.13){$V_{OL}$};
\draw(0,2.4)--(.58,2.4)--(1.12,.78)..controls(1.3,.3)and(1.85,.19)..(2.85,.13);
\draw[dashed](.58,-.1)node[below]{$V_{IL}$}--(.58,2.4);\draw[dashed](1.54,-.1)node[below]{$V_{IH}$}--(1.54,.36);\node[below]at(2.85,0){$V_{DD}$};\draw[dashed](0,.13)--(2.85,.13);

\draw(1.31,.59)--(1.77,.13);\node[right]at(1.66,.9){$a{=}{-}1$};
\end{tikzpicture}


\end{columns}\[V_{OL}\approx\frac12\frac{k_{n,L}}{k_{n,D}}\frac{(V_{DD}-V_{Tn,L})^2}{(V_{DD}-V_{Tn,D})}\]

\end{frame}

\slajdid{09_1-s056}
\begin{frame}[label=09_1-s056]{Uticaj zajedničke osnove}\setlength{\parskip}{2pt}\setlength{\abovedisplayskip}{4pt}\setlength{\belowdisplayskip}{4pt}\setlength{\abovedisplayshortskip}{3pt}\setlength{\belowdisplayshortskip}{3pt}
% element: 09_1-s056:e01
% element: 09_1-s056:e02
% element: 09_1-s056:e03
% element: 09_1-s056:e04
\begin{center}\begin{tikzpicture}[x=.8cm,y=.55cm,font=\sitneoznake,>=Latex,line width=.7pt]\begin{scope}[shift={(0,1.4)}]\draw(-.9,-.2)--(-.15,-.2);\draw(-.15,-.34)--(-.15,.34);\foreach\a/\b in{-.4/-.22,-.1/.1,.22/.4}{\draw(0,\a)--(0,\b);}\draw(0,.3)--(.4,.3)--(.4,.7);\draw(0,-.3)--(.4,-.3)--(.4,-.7);\draw[-{Triangle[length=2mm,width=1.4mm]}](.75,0)--(.03,0);\node[left]at(-.9,-.2){$V_C$};\draw[fill=white](-.9,-.2)circle(.05);\node[right]at(.6,.35){$T_L$};\draw(.75,0)--(.75,-.3)--(.4,-.3);\end{scope}\begin{scope}[shift={(0,0)}]\draw(-.9,-.2)--(-.15,-.2);\draw(-.15,-.34)--(-.15,.34);\foreach\a/\b in{-.4/-.22,-.1/.1,.22/.4}{\draw(0,\a)--(0,\b);}\draw(0,.3)--(.4,.3)--(.4,.7);\draw(0,-.3)--(.4,-.3)--(.4,-.7);\draw[-{Triangle[length=2mm,width=1.4mm]}](.75,0)--(.03,0);\node[left]at(-.9,-.2){$V_I$};\draw[fill=white](-.9,-.2)circle(.05);\node[right]at(.6,.35){$T_D$};\draw(.75,0)--(.75,-.3)--(.4,-.3);\end{scope}\draw(.2,2.1)--(.6,2.1)node[above left]{$V_{DD}$};\draw(.4,-.7)--(.4,-.9)node[ground]{};\draw(.4,.7)--(1.4,.7)node[ocirc]{}node[right]{$V_O$};\draw[DEcrvena](-1.15,-1)--(1.6,2.5);\draw[DEcrvena](-1.15,2.5)--(1.6,-1);\fill(.4,.7)circle(1pt);
\foreach\y in{-.3,1.1}{\fill(.4,\y)circle(1pt);}\end{tikzpicture}
\quad\raisebox{12mm}{$\longrightarrow$}\quad\begin{tikzpicture}[x=.8cm,y=.55cm,font=\sitneoznake,>=Latex,line width=.7pt]\begin{scope}[shift={(0,1.4)}]\draw(-.9,-.2)--(-.15,-.2);\draw(-.15,-.34)--(-.15,.34);\foreach\a/\b in{-.4/-.22,-.1/.1,.22/.4}{\draw(0,\a)--(0,\b);}\draw(0,.3)--(.4,.3)--(.4,.7);\draw(0,-.3)--(.4,-.3)--(.4,-.7);\draw[-{Triangle[length=2mm,width=1.4mm]}](.75,0)--(.03,0);\node[left]at(-.9,-.2){$V_C$};\draw[fill=white](-.9,-.2)circle(.05);\node[right]at(.6,.35){$T_L$};\draw(.75,0)--(.95,0);\end{scope}\begin{scope}[shift={(0,0)}]\draw(-.9,-.2)--(-.15,-.2);\draw(-.15,-.34)--(-.15,.34);\foreach\a/\b in{-.4/-.22,-.1/.1,.22/.4}{\draw(0,\a)--(0,\b);}\draw(0,.3)--(.4,.3)--(.4,.7);\draw(0,-.3)--(.4,-.3)--(.4,-.7);\draw[-{Triangle[length=2mm,width=1.4mm]}](.75,0)--(.03,0);\node[left]at(-.9,-.2){$V_I$};\draw[fill=white](-.9,-.2)circle(.05);\node[right]at(1.02,.25){$T_D$};\draw(.75,0)--(.95,0);\end{scope}\draw(.2,2.1)--(.6,2.1)node[above left]{$V_{DD}$};\draw(.4,-.7)--(.4,-.9)node[ground]{};\draw(.95,1.4)--(.95,-.9)node[ground]{};\draw(.4,.7)--(.83,.7);\draw[preaction={draw=white,line width=2.5pt}](.83,.7)arc[start angle=180,end angle=0,radius=.12];\draw(1.07,.7)--(1.4,.7)node[ocirc]{}node[right]{$V_O$};\fill(.4,.7)circle(1pt);
\fill(.95,0)circle(1pt);\end{tikzpicture}
\end{center}\[V_{Tn,L}=V_{Tn0}+\gamma\left(\sqrt{|2\phi_F+V_{SB,L}|}-\sqrt{|2\phi_F|}\right)\]

\[V_{OH}=V_{DD}-V_{Tn0}-\gamma\left(\sqrt{|2\phi_F+V_{OH}|}-\sqrt{|2\phi_F|}\right)\]
\begin{itemize}\setlength{\itemsep}{1pt}\item $V_{OH}$ značajno opada; slično važi za $V_{OL}$.\item $V_{IL}$ ostaje isti; $V_{IH}$ zavisi od praga opterećenja.\end{itemize}
\end{frame}

\slajdid{09_1-s057}
\begin{frame}[label=09_1-s057]{Obuhvat analize aktivnog opterećenja}\setlength{\parskip}{2pt}\setlength{\abovedisplayskip}{4pt}\setlength{\belowdisplayskip}{4pt}\setlength{\abovedisplayshortskip}{3pt}\setlength{\belowdisplayshortskip}{3pt}
% element: 09_1-s057:e01
% element: 09_1-s057:e02
% element: 09_1-s057:e03
\begin{itemize}
\item Razmotreni su efekti na karakteristiku prenosa.
\item Strujni kapaciteti i dinamički režim nisu računati.
\item Analiza je slična postupku za pasivno opterećenje.
\item Ovim računanjima vraćamo se kod CMOS logičkih kola.
\end{itemize}
\end{frame}

\end{document}
